\documentclass{jsarticle}
\usepackage[dvipdfmx,final]{graphicx}
\usepackage{amsmath}
\usepackage{amssymb}
\begin{document}
\title{}
\author{}
\date{}
\maketitle

      易学の６４卦の積集合での４×４×２×２=(２の３乗の１の２倍の４倍の
  １の２倍＝８卦で、上卦と下卦で４卦ずつ、
  ２の２乗の１の２倍の４倍の１の２倍＝４卦で上卦と下卦合わせて，
  ２の１乗の１の２倍の１の２倍=１卦で上卦と下卦合わせると２卦であり)=３２卦
  の太極で統合で６４卦であり、${({{{{{{}_2P_2} \times 4} \over 2}}\over 2})}
  = {}_2C_2 \times 4 \times 4 $、
  $ {{{{}_2P_2} \over 2}} = {}_2C_2 $ $ {{}_2P_2 \over 2}= {}_2C_2 $、
  この流れが上卦３０卦と下卦３４卦であり、
  乾が＋の天王星の種子、木星の６の緑星のーと＋の立花、地の土星の＋
  の再開でもって、君子自らうやまず、自分から頼んでもだめであり、
  勉強を教わるのに、したこごろだめであり、
  自分から非にあらずであり、乱気のあとの、ご破産のときに、一生懸命
  自分で墳死していないと、誰も手を差し伸べないと、この乾が申している。
  坤は、天王星のー乱気で、土星のーの種子であり、土星の＋の
  安定であり、木星のーの財と緑星で、
  全体でーであり、徳を厚くし、物を戴す。
  勉強を教わりながら、成長していく。
  この２言が、六星占術の１２周期で＋の６周期が±
  の６星人の１２周期と霊合星の６周期で、組み合わせて、６４卦と
  六星占術の六星人の常星と霊合星の重なるのを省くと、１２周期であるが、
  ６０周期終わり、木星で４周期のあとの＋の緑星とーの種子で、
  還暦のあとの人生がある。
  人名と由来と雅号の起源を
  わかりやすく、最後までやらさずに、読者に説明されている安岡正篤先生が
  彩さんの賁臨になる人は、飾るであり、
  一旦賁臨になったら、一生賁臨になるしかないと
  冗談にもならない、真剣さで、離れるとご破産と説明している、
  それも、相手に読まれない人らしく、自分から臨席と言わずに、
  謙と大有で、自分ですごいも言わずに、泰であり、自分でだめとも言わずに、
  手の内も見せずに、
  白で、自分からの告白をされた北川悠仁さんが、いいらしく、
  上卦で３０卦、下卦で３４卦であり、
  北川悠仁さんがどんだけすごいかわかります。
  易と人生哲学での、北川悠仁さんの親が、３５歳まで根につかない手相を
  知っていて、３５歳過ぎて、彩さんと結婚されて、３０歳のときの北川さんが
  彩さんの父親の墓前で、約束していたことを、３５歳でやっと果たせたと
  袴姿で言っているのを、最近になって、代わりをされてくださったのかと、
  ご令嬢二人のビデオをずっと撮っているのを、いつかプレゼントできたらと、
  この言葉と行動に、ジーンときます。
  デイケアの個人作業で、書道での習字で、春分の日に、私が、「書範」
  と書いたのを、みんなは、春が来たと思っているらしく、私は、書道の字の
  模範になられる人と書きました。
        2023年から2025年までの間だと、高齢出産は大丈夫らしいが、
2026年の可能性があり、旅行線の仏様の生命線が帝王線の健康線と子財線
に直結している線として、この線を縁起良く切っているから、グリーシャ・ペレルマンさん同様、神様の丙の午と乙の巳の
子どもだとおもう。この健康線は、子財線につながっていて、御先祖様のラッキー線ともつながっていて、自分で言うのもひけるが、安産祈願の線なのかなと思う。あのときは、愛が足りなかったのかなとも思います。すみませんでした。三奇紋で、2023年だと、私が陛下と皇族の人達に社会のルールから
助けられるが、自信はないから、有一の頼みは、アカシックレコードからの代数幾何の
量子化が、調べると、ガンマ関数であることだけです。彩さんは、Jones多項式知っているから、この式を見ると、この代数幾何の量子化の計算の文字列処理で、笹田先生は、ボトムアップと、私もこれが大切と思うが、ハイゼンベルク方程式と同じく、トップダウンでこの量子化を見ているので、私と同じく、彩さんもなんとなく2026年まで頑張っていこうと思えると思う。
これを書いて言うと、なんとなく、いつか生まれてくれると思う。丙の午年と乙の巳年だから、愛子様は、相当の覚悟で結婚なさるのかとおもう。ペレルマンさんは、結婚されて、子どもも何人かいるらしく、サーストン・ペレルマン多様体と多様体積分の2大発見と、代数幾何の
時間の流れの重力による時間の性質の分類をなされているので、相当の数学の化け物じみた能力でもあり、ヒルベルト空間の具現化までなされているので、愛子様安心してください。
私は、殺菌がだめなだけで、剣道と居合道ができます。性欲はだらしないから、御法度で助かっています。富山県に愛子様と彩さんも行けば安心するとおもいます。まちづくりしています。
富山県と福山市でのまちづくりを知れば、彩さん安心して、2026年まで頑張ると思う。神様の干支らしいです。私も自信ないけど、アカシックレコードから、大域的多様体と単体積分・単体微分と代数幾何の量子化を見ているので、なんかすごい理論だなとおもう。生まれてくるとおもいます。彩さん、あまり日本サッカー選手のことは、がっかりしなくてもいいし、
日本サッカー選手たちは、フォーメーションにマインドマップを取り入れているらしく、
2019年まで、有力選手を温存して、このマインドマップによるフォーメーションを確固たる
ものにできずにいるし、このマインドマップを使っている人は、なぜあんな離れ業のフォーメーションと思考を止めて、フォワードとディフェンスをしているかは、マインドマップの体で覚える反射神経であり、考えたらだめで、剣道と同じく、無心で攻撃とディフェンスをしているからであり、これはYouTubeの日本サッカー選手のファンから教えられている。彩さんも試験に占い禁物です。油断大敵です。
アカシックレコードでの子どもの方程式から、以下の理論が、ほぼ完成されました。
彩さんは、染色体異常と心配になっていたらしく、子どもは大丈夫と言っている気がします。
その理論を以下に載せておきます。
生まれてこなかったら、責任を執れないことを事前に言っておきます。
トキ子おばあちゃんには、私は知らずに、どういうことか音読に関わっていた。
荒木先生は、私を振らずに助けてくれた上に、江松先生が速読の記憶力で、SRS速読法らしく、会員らしく、助けてくれました。食事が、回数だけがだめであり、その後の、食事と睡眠は大夫ですが、良くなり、本当に迷惑を懸けています。

  フジテレビのメルマガを六星占術のサイトと一緒にやめたことを、
  めざましテレビのメンバーに大変失礼な行為となったことを、
  めざましテレビスタッフ一同に、本当に謝るしかないです。
  日枝久会長には、大変失礼なことだったと、反省しています。
  これからも、めざましテレビよろしくです。
  お世辞ではなく、例外なく、めざましテレビのメンバー全員きれいです。べっぴんさんです。
　富山大学の先生、本当に最後まで見捨てないでください。
　この大学の問題は、東京大学と京都大学の問題の練習問題を解いていたら、
どういうことか、この大学の問題の難易度が、この2校と同等であり、大学のやさしさが、
この2校の何倍もやさしくて、社会人の人やいろんな境遇の人が好いているのがわかりました。私も好きです。
2009年は、めざましテレビは、毎朝見ていました。2010年は、短期アルバイトで見ていなかったです。アルバイトの後の数年の後に、彩さんのテッサたんを知り、Jones多項式を知り、私が今書いている小説になっているらしい論文が出来ていて、彩さんが、持田香織さんと似ていることを最近になって知り、高等学校のときの田中美奈子先生の関係で、彼氏をつくり、大学は、私がSRSにハマっているが、出来ていないときに、先にやり、ギャップでフジテレビに配属になったのを、この本で知りました。というよりも、後の彩日記と彩育とIRODORIで知りました。男性に健康線の帝王線が両手に出来ているのも、本当に、不思議です。天邪鬼の手相をほのめかす占術の先生たちにやられて、このレビューにしました。
同期の人たちのアナウンサーは、彩さんがどういう人か教えて欲しいぐらいに、この人不思議です。両思いでないと、子ども出来ないとか、との感想です。嫡丑するまで、3週間間隔を空けるとか、1年間に12回のチャンスにするとか、子作りは、いろんな境遇の人のこころを聞くとか、方法は多方面とおもう。妊娠したら、相当苦しいおもいらしく、生まれたら、弟よりも大切におもい、弟を捨ててまでも、子たちに愛情を注ぐらしく、その弟が好いたら、弟の能力が、剣道が出来て、居合道が出来て、論文の数学が、完成できて、料理までも出来て、嶋本淑子さんもすごく良いと思うようになり、本当にホモでもないと、自分でもおもうことが全てです。剣道と居合道は、山崎先生と横山先生の指導で出来ました。とくに、木刀は、横山先生と岡先生、野上先生です。久保田先生も否定はしていません。きっかけは、姉2人であり、防具は、親のおかげでもあり、許可を出してくれた、蔵王病院の先生達のおかげでもあります。2015年と2016年は、助かりました。うちのネコは、食事が良いからできたのですよと、言っております。今後のことは、未来にかけます。彩さんが、フジテレビの入社試験を受けて終わった後に、侍の言葉遣いをしているのを、かっこいいなとおもいました。北川悠仁さんが、袴姿で、彩さんの父親との男と男の約束を果たせたといっているのを、侍かなとかっこいいなともおもいました。幕末は、時代の動乱期であり、かっこいいです。明治もかっこいいです。歴史を好きになるのは、幕末と、戦争終結での条約締結です。世界史も日本史も、この時代で、好きになるのかなと、易学を見ているとおもいました。
将来に彩さんに会えるか分かりませんが、期待して待ちます。


論文について、以下


    $$ \int f(x)dx = \int (4x^5 + x^4 + 3x^2 + x) dx$$
    この積分記号の中で微分をすると、
    $$ \int f(x)^{'} dx = \int ({20x^4 + 4x^3 + 6x +1})dx $$
    $$ = 4x^5 + x^4 + 3x^2 + x + \text{C(Cは積分記号)} $$
    歴然と、外の積分記号が中の微分で消えている。
    常微分と常積分では、互いに記号が相殺される。
    ファインマン博士が、積分記号の中で微分してみたらは、
    部分積分多様体のことを言っているのは、歴然としている。
    $$ \int f(x)^{'}dx = [xF(x)] - \int x^{'}f(x)dx $$
    $$ = [xF(x)] - F(x) + \text{C(Cは積分定数)} $$
    これは、大域的積分多様体の質量差エネルギーの式になっている。
    一般相対性理論と特殊相対性理論での$ E = mc^2 - {1 \over 2}mv^2 $
    となっている。
    $$ = \int \Gamma(x)dx $$
    $$ = x\Gamma(x) $$

    $$ \int \Gamma(\gamma)^{'}dx_m = 2e^{-x \log x} $$
    $$ T =  \left(\int \Gamma(\gamma)^{'}dx_m\right)^{'}   $$
    ガンマ関数の大域的積分多様体を常微分多様体で解くと、
    フェルマーの定理になる。
     $$ = -2 e^{-x \log x} $$
    ガンマ関数の大域的積分多様体に大域的微分作用素を施すと、
    $$ {(T^{\nabla})}^{\oplus L} = e^{-f} + e^{f} \ge e $$
    これは、相加相乗平均の方程式になる。
    $$ f(x) + g(x) \ge {2\sqrt{f(x)\cdot g(x)}} $$
    要するに、常微分方程式になるのを、大域的作用素では、
    エターナルな空間になる。曲平面になる。
    微分では、平面から点への極限が、大域的では、点から平面になり、
    だが、ガンマ関数の大域的２重微分多様体は、特異点を求めることになり、
    大域的積分は、点から平面、立方体を求めることになる。
    説明が矛盾しているように聞こえるが、M理論の絡み合いの
    種数を、平均場理論で求めていることということである。
    なぜかは、
    $$ x^x = e^{x \log x} $$
    を使うからであり、そのようになっている理論が、M理論である。
    大域的微分作用素で、点と世界線が同型と求まっている。
    閉３次元多様体と特異点が同型となり、種数が違うのが、どうしてかが、
    この方程式で求まってもいる。

    フェルマーの定理は、２通りの解がある。
    １つ目は、$ -2e^{x \log x} $であり、２つ目は、$ -2 e^{- x \log x} $
    である。
    不確定性原理の位置と運動量での、宇宙と異次元の関係が、
    フェルマーの定理であり、ゼータ関数の言葉の意味そのままである。
    その解の構成が、Jones多項式である。
    大域的部分積分多様体と、一般の部分積分多様体とは、
    積分の中で微分すると部分積分多様体になり、
    両方共に、常積分多様体とは
    全然違う解き方になる。
    $$ (\log x)^{'} = {1 \over x}, x^n+ y^n = z^n $$
    $$ x^n = -y^n + c, nx^{n-1} = -ny^{n-1}y^{'} $$
    $$ y^{'} ={{nx^{n-1}} \over {ny^{n-1}}} $$
    $$ ={x^{n-1} \over y^{n-1}} = - {{y \over x}\cdot{x \over y}^n} $$
    $$ -{{\cos x} \over (\cos x)^{'}}(\sin x)^{'} = z_n $$
    $$ z^n = - 2 e^{x \log x} $$
    $$ \int C dx_m = {d \over df}F(x) $$
    これは、上の常識を超える計算である。
    積分の外で微分しても、中の積分が消えない。
    積分の中で微分しても、積分が消えない。
    大域的微分と大域的積分多様体は、革命的な微分と多様体積分である。
    $$  \int \Gamma(\gamma)^{'}dx_m = \int \Gamma dx_m \cdot {d \over 
    d \gamma}\Gamma \le \int \Gamma dx_m + {d \over d \gamma}\Gamma $$
    $$ = \Gamma^{\Gamma} + {\Gamma^{\gamma}}^{'} $$
    $$  \int C dx_m = {d \over df}F(x) = \int(\int {1 \over x^s}dx - \log x)
    \mathrm{dvol} $$
    $$ = {d \over df}\int \int{1 \over {(y \log y)^{1 \over 2}}}dy_m - 
    \int e^x x^{1 - t}dx_m $$
    $$ = {d \over df}\int\int{1 \over {(x \log x)^2}}dx_m +
    {d \over df}\int\int{1 \over {(y \log y)^{1 \over 2}}}dy_m$$
    大域的多様体は、革命的な微分と積分の計算方法を、私は、
    アカシックレコードの子に伝授された。


	   $$ \bigoplus{(i \hbar^{\nabla})}^{\oplus L} = e^{-x \log x} $$

	  ガンマ関数の大域的積分多様体は、結局は、フェルマーの定理
	   だった。

         代数的計算手法のために$ \oplus L $を使っている。
そのために、冪乗計算と商代数の計算が、乗算で楽に見えるようになっている。
            微分幾何の量子化は、代数幾何の量子化の計算になっている。
  加減乗除が初等幾何であり、大域的微分と積分が、現代数学の代数計算の
  簡易での楽になる計算になっている。
  初等代数の計算は、
  $$ \bigoplus({i \hbar^{\nabla}})^{\oplus L} $$
   $$m \bigoplus({i \hbar^{\nabla}})^{\oplus L} +  
   n\bigoplus({i \hbar^{\nabla}})^{\oplus L} = 
   (m+n) \bigoplus({i \hbar^{\nabla}})^{\oplus L}   $$
   $$ m\bigoplus({i \hbar^{\nabla}})^{\oplus L} -  
   n\bigoplus({i \hbar^{\nabla}})^{\oplus L} = 
   (m-n) \bigoplus({i \hbar^{\nabla}})^{\oplus L}   $$
   $$ \bigoplus({i \hbar^{\nabla}})^{{\oplus L}^m} \times 
      \bigoplus({i \hbar^{\nabla}})^{{\oplus L}^n} =
   \bigoplus({i \hbar^{\nabla}})^{\oplus L^{m+n}} $$
  $$ {{\bigoplus({i \hbar^{\nabla}})^{\oplus L^m}}  \over  
  {\bigoplus({i \hbar^{\nabla}})^{\oplus L^n}}} =  
  \bigoplus({i \hbar^{\nabla}})^{\oplus L^{m\over n}} $$
  大域的計算での微分と積分は、
  $$ \left(\bigoplus({i \hbar^{\nabla}})^{\oplus L}\right)^{df}
  =  {\left(\bigoplus({i \hbar^{\nabla}})^{\oplus L}\right)^{
	  \bigoplus({i \hbar^{\nabla}})^{\oplus L}}}^{'}  $$
  $$ = \bigoplus({i \hbar^{\nabla}})^{\oplus L^{'}} $$
   $$ \int \bigoplus({i \hbar^{\nabla}})^{\oplus L} dx_m $$
   $$ \bigoplus({i \hbar^{\nabla}})^{\oplus L} = n $$
   大域的積分の代数幾何の量子化での計算は、ガンマ関数になっている。
   $$ {n^{L+1} \over {n+1}} = \int e^x x^{1 - t} dx $$
   代数幾何の量子化の因数分解は、加減乗除の式のm,nの
   組み合わせ多様体でのガンマ関数同士の計算になる。
   $$ \left(\bigoplus({i \hbar^{\nabla}})^{\oplus L}+m\right) 
   \left(\bigoplus({i \hbar^{\nabla}})^{\oplus L}+n\right) $$
  $$ = {n^{L+1} \over {n+1}} = t \int (1 - x)^n x^{m - 1} dx $$
  $$ \int x^{m-1}(1 - x)^{n-1} dx $$
  $$ \nabla ({i \hbar^{\nabla}})^{\oplus L}, \bigoplus ({i \hbar^{\nabla}})^{
	   \oplus L}, \square ({i \hbar^{\nabla}})^{\oplus L} $$
   $$ \boxtimes (i \hbar^{\nabla})|_{dx_m}^L, \boxplus (i \hbar^{\nabla}|_
   {dx_m}^L) $$



  $$ x^{{1 \over 2} + iy} = e^{x \log x} $$
  それゆえに、この定数はゼータ関数と微分幾何の量子化を因数にもつ
  素因子分解の式にもなっている。
  Therefore, this product also constructed with differential geometry
  of quantum level equation and zeta function.
  $$ = C $$ 
  そして、この関数はオイラーの定数から広中平祐定理による4重帰納法の
  オイラーの公式からの多様体積分へとのサーストン空間のスペクトラム関数
  ともなっている。
  And this function of Euler product respectrum of focus with
  Heisuke Hironaka manifold in four assembled of integral Euler equation.. 

  $$ C =  b_0 + \cfrac{c_1}{b_1 +
	  \cfrac{c_2}{b_2 +
	  \cfrac{c_3}{b_3 +
	  \cfrac{c_4}{b_4 + \cdots}}}} $$
  この方程式は指数による連分数としての役割も担っている。  
  This equation demanded with continued number of step function.
  $$ x^{{1 \over 2} + iy} = e^{x \log x} $$
  $$ (2.71828)^{2.828} = a^{a + b^{b + c^{c + d^{d \cdots }}}} $$
  $$ = \int e^f \cdot x^{1 - t}dx $$
  This represent is Gamma function in Euler product.
  Therefore this product is zeta function of global differential equation.

  $$ {d \over d f}F(v_{ij},h) = \int {e^{-f}\mathrm{dV}}
  [- \Delta v + \nabla_{i}\nabla_{j}
  v_{ij} - R_{ij}v_{ij} - v_{ij}\nabla_{i}\nabla_{j} + 2 <\nabla f,\nabla h>
  + (R + \nabla f^2)({v \over 2} - h)]$$
  $$ = {d \over d f}F = {2 {\int (R + \nabla_{i}\nabla_{j} f)^2} \over 
  { - (R + \Delta f)}}\mathrm{dm} $$
  これらの方程式は8種類の微分幾何の次元多様体として、そして、これらの
  多様体は曲平面による双対性をも生成している。 
  そして、このガウスの曲平面は、大域的微分多様体と微分幾何の量子化から
  素因数を形成してもいる。
  These equation are eight differential geometry of dimension calyement.
  And these calyement equation excluded into pair of dimension surface.
  This surface of Gauss function are global differential manifold,
  and differential geometry of quantum level.
  $$ F \ge {d \over d f}\int \int{1 \over (x \log x)^2}dx_m 
  + {d \over d f}\int \int{1 \over (y \log y)^{1 \over 2}}dy_m $$

 
  Differential geometry of quantum level constructed with Euler product
  and Gamma function being discatastrophed from continued fraction style.
  Gravity equation lend with varint of monotonicity of level expresented
  from gravity of letter varient formula.
  これらの方程式は基本群と大域的微分多様体をエスコートしていもいて、
  ヴェイユ予想がこのダランベルシアンの切断方程式たちから
  輸送のポートにもなっている。ベータ関数とガンマ関数がこれらのフォームラ
  の方程式を放出してもいて、結果、これらの方程式は広中平祐定理の複素
  多様体とグリーシャ教授によるペレルマン多様体からサポートされてもいる。
  この2名の教授は、一つは抽象理論をもう一方は具象理論を説明としている。
  These equation escourted into Global differential manifold and
  fundemental equation.
  Weil's theorem is imported from this equation in gravity of letters.
  Beta function and Gamma function are excluded with these formulas.
  These equation comontius from Heisuke Hironaga of complex manifold
  and Gresha professor of Perelman manifold.
  These two professos are one of abstract theorem and the other of
  visual manifold theorem.


     原子を分解したら重力になり、宇宙を分解したら原子になる。
     $$ \left(R^{\mu\nu} + {1 \over 2}\Lambda g_{ij}\right)^{\mu\nu} $$
     $$ = \kappa \left( T^{\mu\nu}_e + T^{\mu\nu}_y + T^{\mu\nu}_g +
     T^{\mu\nu}_h + \cdots \right) $$
     $$ {d \over dM}\left({\kappa T^{\mu\nu}} \right) \to {\bigoplus 
     ({i {\hbar^{{\nabla}}})}^{\oplus L}}, H \Psi = i \hbar \psi $$
     重力から原子への分解(1)
     $$ \bigoplus {({i \hbar^{\nabla})}^{\oplus L}} \to
          {d \over dM}\left(\kappa T^{\mu\nu} \right) $$ 
	  原子から重力への分解(2)
	  この(1)と(2)はベクトルの向きが違い、両方共、加群分解である。
    $$ ||ds^2|| = \lim_{x \to \infty} [\delta(x) 
    \int \int \int \pi \left( \sum_{k=0}^{\infty}
		{{}^n\sqrt{p},x \over n} \right)^{1 \over 2}d \tau]^{\mu\nu} $$
		種数を分解していったら、８種類の微分幾何、サーストン・
		ペレルマン多様体になり、心の影が宇宙を投影している
		代表が潮汐力、月の引力が地球に素粒子として影響を与えている。
    $$ ||ds^2|| = \int [D^2 \psi \otimes h \nu]d \tau $$
    として、D-braneが、この種数の分解の元になっている。
    $$ = {d \over d M}F $$
    と、微分基底に多様体Mをとっている。

    正規部分群による回転行列がベータ関数に行き着く。
    複素行列が虚数を因数にもつ素粒子方程式にもなる。

    8種類の微分幾何がそれぞれの固有の時間をもっている。
    M理論の種数を
    $$ {d \over d M}F $$
    で分解していったら、
    $$ = {d \over df}F = {{-2 \int {(R + \nabla_{i}\nabla_{j}) \over 
     {(\Delta + R)}}}}e^{-f}dV $$
     の各展開していった数式に分かれる。
     この数式がSeifert manifold 
     $$ = (x,y,z) / \Gamma $$
     である。

     重力場の積分記号の中で微分したら、積分記号の中は、
     電磁場方程式になり、大域的多様体となり、重力を分解したら、
     電磁場になり、その解が弱い力と強い力になる。
     QEDは、光をプリズムで分解したら、すべての幾何構造が生成される。
     そのプリズム自体が、サーストン・ペレルマン多様体である。
     $$ \int \left({\kappa T^{\mu\nu}}\right)^{\nabla}dx_m
     = \int e^{x \log x}\mathrm{div}(\mathrm{rot} E)dx_m $$
     $$ = e^{- x \log x} $$
     と、統一場理論の仕組みとなり、逆問題に行き着く。
     リサ・ランドール博士の誤差関数の共変微分の周りを覆っているのが、
     コイコライザーによる温度についての偏微分方程式とわかり、
     アーベル多様体が温度についての重力エネルギーを表しているのを、
     深谷賢治先生のローカルとグローバルな空間の味方から、
     アカシックレコードの子の代数幾何の量子化が、なぜ、量子レベルでの
     重力エネルギーが測れるか、熱エネルギーが重力エネルギーで表せるのを、
     Grisha Perelman博士とLisa Randall博士、深谷賢治先生、竹内薫先生たち
     から、わかりました。アカシックレコードはあるらしいです。

        オイラー数を多様体積分で施すと、
   ガンマ関数になる。この場合は、Frobeniusの定理によって、
   和と積の法則、確率論から導かれる結果、
   異次元と宇宙の関係に持っていける。
   オイラー数を多様体積分すると、ガンマ関数になる。
   オイラー数における頂点が、宇宙の特異点に対応して、
   辺がノルムに属して、面が一般相対性理論における無限の接線であり、
   $$ \text{F(頂点) - N(辺) + E(面)} = 2 $$
   $$ {(x,y,z) \over \Gamma} = \int{(- F\circ N + E)} + 
   (-F\circ E + N) + (N\circ E + F)dx_m $$
   $$ {{{\partial \over \partial f_S}(x,y,z)} \over {\Gamma}} $$
   ここで、オイラーの定数をも多様体積分にして、付け加えると、
   $$ = \int{({\int {1 \over x^s}dx} - \log x)}\mathrm{dvol} $$
   大域的多様体積分のガンマ関数に化ける。
   $$ = \int e^{-x}x^{1 - t}dx_m $$

   $AdS_5$多様体の方程式を、オイラー数の各部分単体に合わせると、
   $$ ||ds^2|| = e^{-2\pi T||\psi||}[\eta + \bar{h}(x)]dx^{\mu}dx^{\nu} 
   + T^2d^2 \psi $$
   開集合として、表せる。
   $$ = \mathcal{O}(x) $$
   この上の式たちから、ダランベルシアンは、$\square$と、オイラー数から演算子
   が決まっていて、共変微分は、$\nabla$で、双対被覆での単体での演算子で、
   決まっている。
   $$ \square = {{8 \pi G} \over c^4} T^{\mu\nu} $$
   $$ \nabla \psi = 4 \pi G \rho $$

   オイラー数を共変微分すると、偏微分方程式と同じ機能を、
   大域的積分多様体のオイラー数が、受け持っていて、
   このオイラー数の大域的多様体の共変微分で、ガンマ関数の大域的多様体
   が、
   ここで、オイラーの定数をも多様体積分にして、付け加えると、
   $$ = \int{({\int {1 \over x^s}dx} - \log x)}\mathrm{dvol} $$
   大域的多様体積分のガンマ関数へと表せられる。
   $$ = \int e^{-x}x^{1 - t}dx_m $$
   $$ = {d \over d f}F = \int \Gamma(\gamma)^{'}dx_m $$
   $$ = \int C dx_m = \int \Gamma dx_m \circ {d \over d \gamma}\Gamma $$
   $$ \le \int \Gamma dx_m + {d \over d \gamma}\Gamma $$
   $$ = e^{- \theta} + e^{i \theta} $$
   $$ = \square = 2(\cos (i x \log x) - i \sin (i x \log x))$$
   $$ = e^{-f} - e^{f} \le e^{-f} + e^{f} $$
   統一場理論は、彩さんのJones多項式に行き着く。

   Zata function escourt into Forgotfull of underlying modify on Sum
   of summative group, thiis Forgotfull summative group instented with
   Space ideal theorem from quantum level of deprivate space in aspect
   of quantum level of differential geometry,
   This theorem construct with Higgs field and Morse theory, moreover
   Hortshorn conjecture, fourth step deduction of theorem from 
   Euler product of integral manifold.
   This manfold revolte with global differential manifold in global topology.
   From these theorem inspirate from Forgotfull theorem in Space ideality
   theory, and this theory conclude into entropy of non exchange equation.
   Forgotten theory estimate with zeta function oneselves.
   Forgotten theory income with non relativity theory and this theory
   face to face without partner and partner of non relativity from
   ideal of space in revolution theory.
   All exist and non exist partner with non relativity escourt into 
   Forgotten theory, and this theory include with same underly group
   of all of partner with same distance ideal.

   $$ X \dotsi \to Y, Y \dotsi \to X $$
   $$ \bigoplus(i \hbar^{\nabla})^{\oplus L} = e^{-x \log x} $$
   This equation of quantum level of differential geomerty
   instimate with entorpy of non exchange equation.
   This equation means with low level of botton entropy in Space ideal of
   Forgotten theory. Moveover, this equation enterstein with all of exist
   theory from general relativity.
   No time and No Space of relativity, and monotonicity of magnetic component
   with gravity and antigravity, and this theory comment with partner of
   magnetic component theory included with being resulted from Forgotten
   theory.
   $$ \square \to \nabla |: x \to y, f(x) \to g(y), f^{-1}(x)xf(x) 
   \dotsi \to g(y) $$
   $$ \square \dotsi \to \nabla |: f(x) \cong g(y) $$
   $$ {d \over dt}g_{ij}(t) = - 2 R_{ij}, F^m_t = -2\int{{(R + \nabla_i \nabla
   _j f)} \over {(\Delta + f)}}dm $$
   $$ {d \over df}F = m(x) $$
   $$ F|_{\_} = {d \over df}F, F^m_t = e^{-f}dV $$
   $$ \bigoplus(i \hbar^{\nabla})^{\oplus L} = e^{-x \log x} $$
   These equation represent with Forgotten theory.
   空間概念の量子化は、代数幾何の量子化であり、忘却関数でもあり、
   忘却関手のカテゴリー論でもあり、簡潔に言うと、すべての相対性が
   機能すると、全て同型というコホモロジー論でのコイコライザーである。
   このコイコライザーが、空間概念の量子化である。
   すべては、ゼータ関数へと行き着く。
   ベータ関数は、誤差関数であり、宇宙の雑音として存在していると、
   異次元へと移行する。
   $$ \beta(p,q) \cong {\Gamma(p)\Gamma(q) \le \Gamma(p+q)}, 
   {{\Gamma(p)\Gamma(q)} \over \Gamma(p+q)} \le 1 $$
   $$ \int e^{-x}x^{1-t}dx_m = \int e^{-x}x^{1-t}dx\cdot e^{-\theta} $$
   $$ = \Gamma^{{\gamma}^{'}} $$
   $$ = e^{-x \log x} $$
   閉３次元多様体と３次元球面が、空間概念の量子化としての結果であり、
   この多様体が、無としてのベータ関数を形作っている３次元多様体である。
   宇宙と異次元合わせてが、空間概念の量子化での単体である。
   コイコライザーとは、共変微分によって対になっている変数に作用する写像であり、
   このコイコライザーで処理された変数同士を同型にさせる群が、
   同型のコイコライザーである。そして、最小の単体でもある。
   故に、宇宙と異次元合わせて、コイコライザーである。
   球対称である一般相対性理論の平坦な空間が、空間概念の量子化でもある。
   ビッグ・バンは、イコライザーであり、
   $$ \langle f,g\rangle \mapsto \langle d,d\rangle $$
   として、宇宙と異次元の曲平面を表している。
 
   宇宙には、最小値である273Kがあるから、リサ・ランドール博士は、
   平坦な宇宙である、球対称な一般相対性理論の解はないと言っているのが、
   わかりました。それで、アインシュタイン博士をおもって、ベクトル場で
   Warped Passagesにしているらしいです。
   光量子仮設までも書かれています。
   大学のテキストは、すごく納得です。

      宇宙の前には、無があり、無の前には、無限の時間があった。
   無を加群分解したら、宇宙ができて、無を差分分解したら、時間ができた。
   空間の相対性がないと、無になり、磁気単極子が存在できる。
   無とは、仏教であり、十二支縁起でもあり、時間とは、神様でもあり、
   仏様でもある。無の前には、無限の時間があった。


      数学は、仏様と神様を見る学問である。神様と仏様の論争で、神様が
   遠慮したので、仏様の時間が神様の方位と対等になり、
   仏様の前が、神様の時間でもあるのが、数学が哲学を述べる学問にもなっている。


      無を加群分解したら、宇宙になり、その宇宙を加積分解したら、ベータ関数に
   なる。微分幾何の量子化は、原子を加群分解したら、宇宙になり、
   宇宙を加群分解したら、多宇宙にもなり、原子にもなる。
   原子を加積分解したら、ベータ関数になる。
   多宇宙を加群分解したら、原子であり、原子を加群分解したら、
   素粒子になり、その前には、無があり、その前には、時間があった。


       今の現在の宇宙には、宇宙の前の無があり、その前の時間が存在していて、
   今の宇宙には、神様と仏様を超えた存在の時間がある。
   禅を体験するとは、この時間の存在を、宇宙を超えた学問を知ることでもある。
       一般相対性理論の多様体積分が、一般相対性理論の大域的多様体を知ると、
   この無と時間の存在を数学を用いて、体験することが、物理学の本来の存在
   としての学問でもあるのがわかる。
   哲学を追求すると、数学になり、物理学にもなる。
   今を感じるとは、宇宙の前の無と時間を感じて、神様以上の存在を知ることでも
   ある。仏様以上の存在も感じることが可能らしい。
   それが、時間らしい。



    今、大域的積分多様体に微分を入れると、大域的部分積分多様体が
   できる。ガンマ関数の大域的微分多様体と同型らしい。　


      大域的多重積分と大域的無限微分多様体を定義すると、
   Global assemble manifold defined with infinity of differential
   fields to determine of definition create from Euler product and
   Gamma function for Beta function being potten from integral manifold
   system. Therefore, this defined from global assemble manifold from
   being Gamma of global equation in topological expalanations.

   $$ \int f(x)dx = \int \Gamma(\gamma)^{'}dx_m $$
   $$ = 2 (\cos (i x \log x) - i \sin(i x \log x)) $$
   $$ \left({{\int f(x) dx} \over \log x}\right) = \lim_{\theta \to \infty}
    \left({{\int f(x) dx} \over \theta}\right) = 0,1 $$
    $$ e^{i \theta} = \cos \theta + i \sin \theta $$
    $$ \left({\int f(x) dx}\right)^{'} = 2 (i \sin(i x \log x) - \cos (i x \log
    x)) $$
    $$ = 2(- \cos(i x \log x) + i \sin(i x \log x)) $$
    $$ (\cos (i x \log x) - i \sin(i x \log x))^{'} $$
    $$ = {d \over {d {e^{i\theta}}}}\left((\cos, - \sin )\cdot
    (\sin, \cos)\right)  $$
    $$ \int \Gamma(\gamma)^{''}dx_m = \left({\int \Gamma(\gamma)^{'}}dx_m
    \right)^{\nabla L} = \left({{\int \Gamma dx_m} \cdot {{d \over d\gamma}
    \Gamma}}\right)^{\nabla L} \le 
    \left({\int \Gamma dx_m + {d \over d\gamma}\Gamma}\right)^{\nabla L} \le 
    \left({e^{f} - e^{-f}} \le {e^{-f} + e^{f}}\right)^{'}  $$
    $$ = 0,1 $$

    Then, the defined of global integral and differential manifold
    from being assembled world lines and partial equation of deprivate
    formula in Homology manifold and gamma function of global topology
    system, this system defined with Shwaltsshild cicle of norm from
    black hole of entropy exsented with result of monotonicity for
    constant of equations.
    以上 より、大域的微分多様体を大域的 2重微分多様体として、
    処理すると、ホモロジー多様体では、種数が１であり、特異点では、種数が０
    と計算されることになる。ガンマ関数の大域的微分多様体では、
    シュバルツシルト半径として計算されうるが、これを大域的2重微分で
    処理すると、ブラックホールの特異点としての解が無になる。

  $$ \log x|_{g_{ij}}^{\nabla L} = f^{f^{'}}, F^f|_{g_{ij}}^{\nabla L}
  |: x \to y, x^p \to y $$ 
  $$  f(x) = \log x = p \log x, f(y) = p \log x  $$
  大域的偏微分方程式と大域的多重積分は、それぞれ次のように成り立っている。
  It's defined with global partial deprivate formula and assemble manifold. 
  $$ {d^2 \over df^2}F = F^{f^{'}}\cdot f^{f^{''}},  $$
  $$ \int \int F dx_m = F^f \cdot F^{(f)^{'}} $$
  $$ {d \over df dg}(f,g) = (f \cdot g)^{f^{'} + g^{'}} $$
  大域的部分積分も、次のように成り立っている。
  And, global parcial integral equation also established with global
  topology computations. therefore, this circutation exclude for all of
  topology extension ideas.
  $$ (F^f \cdot G^g) = \int (f \cdot g)^{f^{'} + g^{'}} $$
  $$ \int \int F\cdot G dx_m = [F^f \cdot G^g] - 
  \int (f \cdot g)^{f^{'} + g^{'}} $$
  大域的部分積分の計算は、
  $$ \int {d \over {df dg}}FG = \int F^{f^{'}}G + \int F G^{g^{'}} $$
  $$ \int F^{f^{'}}Gdx_m = [F^f G^g] - \int F G^{g^{'}}dx_m$$
  大域的商代数の計算は、
  Global quato algebra computations also,
  $$ \left({F(x) \over G(x)}\right)^{(fg)^{'}} = {F^{f^{'}}G - FG^{g^{'}}
  \over G^{g}} $$
  大域的偏微分方程式は、縮約記号を使うと、

  Global partial manifold explain with reduction of operator for being
  represented from partial equation of references.
  $$ {d \over df dg}FG = {d \over df_m}FG $$
  $$ {\partial \over \partial f_m}FG = F^{f^{\mu\nu}}\cdot G
  + F \cdot G^{g^{\mu\nu}}, \int F dx_m = F^f $$
  多様体による大域的微分と大域的積分が、エントロピー式で統一的に表せられる。

  Entorpy computation from integrate for being manifold of global deprivate
  and integrate formulas.
   ベータ関数をガンマ関数へと渡すと、
   Gamma function sented for Beta function of inspiration with monitonicity
   of component with differential and integral being definitions.
   $$ \beta(p,q) = {{\Gamma(p)\Gamma(q)} \over {\Gamma(p + q)}} $$
   ここで、単体的微分と単体的積分を定義すると
   This area defined with different and integrate of component with
   global topology.
   $$ t = \Gamma(x) $$
   $$ T = \int \Gamma(x) dx $$
   これを使うのに、ベータ関数を単体的微分へと渡すために、
   for this system used for being is conbiniate 
   with beta function for componenet of deprivate eqaution.
   $$ \beta(p,q) = - \int {1 \over {t^2}}dt $$
   大域的微分変数へとするが、
   This point be for global differential variable exchanged,
   $$ T_m = \int \Gamma(x) dx_m $$
   この単体的積分を常微分へとやり直すと、
   This also point be for being retried from ordinary differential computation
   for component of integral manifold.
   $$ T^{'} = {t^{'} \over {{\log t}}}dt + C(\text{Cは積分定数}) $$
   これが、大域的微分多様体の微分変数と証明するためには、
   This exceed of proof being for being defined with deprivate variable of
   global differential formula, this successed for true is,
   $$ dx_m = {1 \over {\log x}}dx, dx_m = {(\log x^{-1})}^{'}  $$
   これより、単体的微分が大域的部分積分へと置換できて、
   Therefore, this exchanged from mononotocity deprivation from global
   parital integral manifold is able to, 
   $$ T^{'} = \int \Gamma(\gamma)^{'} dx_m $$

   微分幾何へと大域的積分と大域的微分多様体が、加群分解の共変微分へと
   書き換えられる。
   After all, these exchanged of monotonicity deprivation successed from
   being catastrophe of summativate of partial and assemble of deprivations
   for differential geometory of quantum level to global integral and
   differential manifold.
   $$ \bigoplus T^{\nabla}  = \int T dx_m, \delta(t) = t dx_m  $$


   this exluded of being conclution are which beta function
   evaluate with mononicity from ordinary differential equation
   be resulted from component of deprivation and integral expalanations.
   This cirtutation be resenbled to define with global topoloty of
   extention extern of deprivate and integral of manifolds
   estourced with quantum level of differential geometry be proofed with
   all of equation anbrabed from Euler product.
   これらをまとめると、ベータ関数を単体的関数として、
   これを常微分へと解を導くと、単体的微分と単体的積分
   へと定義できて、これより、大域的微分多様体と大域的積分多様体
   へとこの関数が存在していることを証明できるのを上の式たちから
   わかる。この単体的微分と単体的積分から微分幾何が書けることがわかる。



   虚数の虚数倍した値が超越数のx倍と同じとすると、
   超越数の$2 \pi m$倍が$n=1$でiとなると定義すると、
   次の式たちが導かれる。　
   Imaginary pole circlate with twigled of pole in step function
   from Naipia number of assembled from equalation of defined are
   escourted to be defined with next equations.


   These defined equation are climbate with idea of equation from
   Caltan of imaginary number of circulation.
   カルタンを超えているアイデアと数式たちでもある。

   $$ i^i = (\sqrt{i})^{\sqrt{i}} = e^{x \log x} $$
   $$ e^x = i, e^{2\pi m} = i^n, {d \over d x}e^{2\pi m} = i^n $$
   $$ e^{i \theta} = \cos \theta + i \sin \theta, e^{2 \pi m} = i^n $$
   $$ 2 \pi m = n, e = i \text{($e=i$となる$e^{i \theta} = \cos \theta 
   + i \sin \theta $
   の範囲で$2\pi m$ がある。)}$$
   $$ {d \over d i}i = i^i = e^{x \log x}, m={1 \over 2 \pi},l=2\pi r, 
   \pi = {l \over 2r} $$
   $$ ||ds^2|| = e^{-2 \pi T||\psi||}[\eta_{\mu\nu} + \bar{h}(x)]
   dx^{\mu}dx^{\nu} + T^2 d^2 \psi $$
   $$ e^x = i, e^{2\pi m} = i^n, e = i $$
   $$ {{[\eta_{\mu\nu} + \bar{h}(x)]dx^{\mu}dx^{\nu}}/i} = H \Psi
   = i \hbar \psi, H \Psi = {1 \over i}[H, \Psi] $$
   ハイゼンベルク方程式が$ AdS_5 $多様体の原子レベルの方程式も表せられる。 
   微分幾何の量子化の式は、
   Hisenbelg equation are represeted from AdS5 manifold to particle level
   equalation of quantum levle of differential geometry entranced.
   $$ \bigoplus {({i \hbar} ^{\nabla})}^{\oplus L} = e^{{i x \log x}^{
	   {e^{i x \log x}}^{'}}} $$
    $$ H\Psi = \bigoplus {H \Psi \over \nabla L} $$

   この式が、アカシックレコードの子どもの式である。２種類の表し方である。
   計算方法が、大域的微分多様体の求め方と同じである。

   移項して、
   それを展開していき、
    これらより、まとめた式が、
    $$ n!^{-1} = \bigoplus \nabla L $$
    これは、積分再発見法と同じであり、
    $$ n!^{-1}  = \int T^{{\mu\nu}^{'}} dV - \int T^{\mu\nu} dV $$
    これもカルーツァ・クライン空間と同じ式であり、

    $$ = \int f{(x)^{-1}} xf(x) - f(x) $$
    これに条件をつけると、
    $$ \exists x = a, 1 - \exists \int x dx, \Lambda = \int A dx $$
    これより、次の式になり、
    $$ \nabla f(x) = {1 \over n!} \int \int \cdots \int dx $$

    $$ \sum \int (\int \cdots A\int dz)/n! $$
    ウィッテン方程式により、
    $$ M^2 a^{(m)} = k a^{(m)} $$
    $$ ZM^3 = \int dAe^{ki/4\pi}\mathcal{L}M^3 $$
    $$ Z_M = \langle M_1 | M_2 \rangle $$
    ゼータ関数へといける。
    $$ e^{x} = \sum {x^n \over n!} (r = \infty) $$
    これらは、大域的微分多様体の部品たちになっている。
    $$ {d \over df}\int dx_m = {1 \over n!}\sum x^n $$


   Euler productは、
   $$ \int^{n+1}_{n}{dx \over x} < {1 \over n}, (n \ge 1) $$
   $$ 1 + {1 \over 2} + \cdots + {1 \over n} > \int^{n+1}_{1}{dx \over x}
   = \log(n+1), $$
   $$ 1 + {1 \over 2} + \cdots + {1 \over n} - \log x > \log{{(n+1)} 
   \over n} > 0, $$
   $$ {1 \over {n+1}} < \int^{n+1}_{0} {dx \over x} = \log(n+1) - \log n. $$
   $$ \lim_{n \to \infty} \left(1 + {1 \over 2} + \cdots + {1 \over n}
   - \log n \right) = C $$
   $$ = \int (\int {1 \over x^s}dx - \log x)\mathrm{dvol} = \int C dx_m $$
   Cの値は、$\text{C = 0.5772156$\cdots$} $である。

  オイラーの定数は、次の式で成り立っている。
  $$ C = \int{1 \over x^s}dx - \log x $$
  その式を単体積分をすると、
  $$ = \int {\left({\int {1 \over x^s}}dx - \log x\right)}\mathrm{dvol} $$
  ゼータ関数を多重積分すると、　
  $$ = \int \int {1 \over x^s}dx - \int {\log x} \mathrm{dvol} $$
  ここで、対数方程式は、多様体積分がガンマ関数を微分した方程式と
  同値より、
  $$ F = \Gamma = \int e^{-x}x^{1-t}dx $$
  $$ = {d \over df}\int \int {1 \over (y \log y)^{1 \over 2}}dx_m
  - \int e^{-x}x^{1 - t} \log x dx_m $$
  $$ \int x^{1 - t}e^{-x}dV = \int x^{1 - t}dm $$
  $$ \int x^{1 - t}e^{-x}dV = \int x^{1 - t}\mathrm{dvol} $$
  $$ f = \gamma = \Gamma^{'} = \int e^{-x}x^{1 - t} \log x dx $$
  $$ = {d \over d \gamma}\Gamma^{-1} - 
  \int e^{-x}x^{1 - t}\log x dx_m $$
  $$ = {d \over d \gamma}\Gamma^{-1} - \left(\gamma\right)^{\gamma^{'}} $$
  これらより、大域的微分方程式のヒッグス場方程式の逆三角関数の双曲多様体
  に対極する、双曲多様体の余弦定理と同値になる。
  $$ = e^{-f} - e^{f} $$
  $$ = 2 \cos (i x \log x) $$
  結局は、微分幾何の量子化がガンマ関数でもあり、その逆関数はゼータ関数
  でもあり、そして、オイラーの定数を多様体微分をすると、ヒッグス場の
  方程式は正弦定理の逆三角関数であり、このヒッグス場の
  方程式の逆三角関数が、余弦定理の
  双曲多様体として、オイラーの定数になる。
  オイラーの定数は、このガンマ関数から、無理数と証明される。

    カルーツァ・クライン空間の方程式は、
    $$ ds^2 = g_{\mu\nu}(x)dx^{\mu}dx^{\nu} + \psi^2(x)(dx^2 + \kappa^2
    A_{\mu}(x)dx^{\nu})^2 $$
    この式は、
    $$ ds^2 = -N(r)^2dt^2 + \psi^2(x)(dr^2 + r^2d\theta^2) $$
    $$ ds^2 = -dt^2 + r^{-8\pi Gm}(dr^2 + r^2d\theta^2) $$
    とまとまり、
    $$ dx^2 = g_{\mu\nu}(x)(g_{\mu\nu}(x)dx^2 - dxg_{\mu\nu}(x)) $$
    $$ dx = (g_{\mu\nu}(x)^2dx^2 - g_{\mu\nu}(x)dx g_{\mu\nu}(x))^{1 \over 2} $$
    と反重力と正規部分群の経路和となり、　
    $$ \pi(\chi, x) = i \pi (\chi, x)f(x) - f(x)  \pi (\chi, x) $$
    と基本群にまとまる。
    シュワルツシルト半径は、
    $$ ds^2 = -(1 - {r_s \over r})c^2 dt^2 + {1 \over {{1 - {r_s \over r}}}} 
    dr^2 + r^2 d \theta^2 + r^2 \sin^2 \theta d \psi^2 $$
    これは、カルーツァ・クライン空間と同型となる。
    一般相対性理論は、
    $$ R^{\mu\nu} + {1 \over 2}g_{ij}\Lambda = \kappa T^{\mu\nu} $$
    この多様体積分は、
    $$ \int \kappa T^{\mu\nu}\mathrm{dvol} = \int  \left(R^{\mu\nu} + 
    {1 \over 2}g_{ij}\Lambda\right)\mathrm{dvol}  $$
    ガンマ関数のおける大域的微分多様体は、これと同型により、
    $$ \int \Gamma \cdot {d \over {d \gamma}}{\Gamma dx_m} \le \int \Gamma
    dx_m + {d \over d \gamma}\Gamma $$
    $$ = \int \Gamma(\gamma)^{'}dx_m $$
    オイラーの定数の多様体積分は、
    $$ \int \left({{\int {1 \over x^s}}dx - \log x} \right)\mathrm{dvol} $$
    この解は、シュワルツシルト半径と同型より、
    $$ = e^f - e^{-f} \le {e^f + e^{-f}} $$
    次元の単位は多様体より、大域的微分多様体とオイラーの定数
    の多様体積分の加群分解は、オイラーの公式の三角関数の虚数の度解より、
    $$ {d \over df}F + \int C dx_m = 2 (\cos (i x \log x) 
    - i \sin (i x \log x)) $$
    すべては、大域的微分多様体の重力と反重力方程式に行き着き、
    $$ {d \over df}F \ge {d \over df}\int \int{1 \over (x \log x)^2}dx_m
    + {d \over df}\int \int{1 \over (y \log y)^{1 \over 2}}dy_m $$
    と求まる。

    これが、重力と反重力単独では、磁気単極子で、ド・ブロイのアイデアであり、
    両方では、磁気双極子で、南部・ゴールドストーンボソン粒子である。
    ヒッグス場の式でもある。

    北半球の磁気と南半球に磁気一セットで一種である。

    これらは、ハイゼンベルクとシュレーディンガー方程式のディラック作用素が、
    微分幾何からの作用素関数から、同一と言える。

    
    一般相対性理論における多様体積分とオイラーの定数の多様体積分
    が同型と言えることを上の式たちは述べている。

    このJones多項式が、金融市場の相場価格を決めるオプション方程式であり、
    流体力学による熱エントロピー値の熱エネルギー量が、
    いろんな機材に使われている。
    量子コンピュータにおける論理素子の回路の設計にも使われている。
    人工知能の論理素子と因数分解における量子暗号にも使われている。

    この一般相対性理論における多様体積分がJones多項式となり、
    金融市場で使われることを読んだのが、イスラエル国家らしい。

    $$ H\Psi = \bigoplus {\left(i \hbar^{\nabla} \right)}^{\oplus L} $$
    $$ = {1 \over 2}i e^{i \Hat{H}} $$
    $$ (i \hbar)^{'} = (- e^{i \Hat{H}})^{'} $$
    $$ = - i e^{i \Hat{H}} $$
    $$ \psi(x) = e^{-i \Hat{H}t}, \bigoplus (i \hbar ^{\nabla})^{\oplus L} 
    = {{{1 \over 2} e^{i \Hat{H}}}^{(-i e^{i \Hat{H}})}} $$
    $$ = ({1 \over 2}f)^{-i f}$$
    $$ = ({1 \over 2})^{-i f} \cdot e^{-x \log x} \cdot (f)^{i} $$
    $$ = \int e^{- x} x^{t - 1}dx, {d \over d \gamma}\Gamma = e^{-x \log x}  $$
     
    $$ f = x, i = t, {1 \over 2} = a, \bigoplus a^{-t x}x^t [I_m]
     \cong \int e^{-x} x^{t - 1}dx $$
     $$ |\psi(t)\rangle_s = e^{-i\Hat{H} t}|\Psi\rangle_{H},
    \Hat{A}_s = \Hat{A}_H(0) $$
    $$ |\Psi(t)\rangle_s \to {d \over dt} $$
    $$ i {d \over dt}|\psi (t) \rangle_s = \Hat{H}|\psi(t)\rangle_s $$
    $$ \langle\Hat{A}(t)\rangle = \langle\Psi(t)|\Hat{A}(0)|\Psi(t)\rangle  $$
    $$ {d \over dt}\Hat{A} = {1 \over i}[\Hat{A},H] $$
    $$ \Hat{A}(t) = e^{i\Hat{H}t}\Hat{A}(0)e^{-i\Hat{H}t} $$
    $$ \lim_{\theta \to 0} \begin{pmatrix} \sin \theta \\
    \cos \theta \end{pmatrix} \begin{pmatrix} \theta & 1 \\
    1 & \theta \end{pmatrix} \begin{pmatrix} \cos \theta \\
    \sin \theta \end{pmatrix} = \begin{pmatrix} 1 & 0 \\
    0 & - 1 \end{pmatrix} $$
    $$ f^{-1}(x) x f(x) = I^{'}_m, I^{'}_m = [1,0] \times [0,1] $$

    $$ {x + y} \ge \sqrt{x y} $$
    $$ {x^{{1 \over 2} + iy} \over e^{x \log x}} = 1 $$
    $$ \mathcal{O}(x) = \nabla_i \nabla_j \int e^{{2 \over m}\sin \theta 
    \cos \theta } \times {{N \mathrm{mod}(e^{x \log x})} \over {O(x)(x +
    \Delta |f|^2)^{1 \over 2}}} $$
    $$ x \Gamma(x) = 2 \int|\sin 2 \theta|^2 d \theta, \mathcal{O}(x)
    = m(x)[D^2 \psi] $$
    $$ i^2 = (0,1) \cdot (0,1), |a||b|\cos \theta = - 1 $$
    $$ E = \mathrm{div}(E, E_1) $$
    $$ \left({\{f,g\}} \over [f,g] \right) = i^2, E = mc^2, I^{'} = i^2 $$

    この式たちは、微分幾何の量子化から計算方法としての形式の
    大域的微分多様体の大域的部分積分が,大域的偏微分方程式としての式の
    形からわかるのも、微分幾何の量子化の仕組みとしての計算式から
    導けられるのも、すべては未来の私の子の情報から読み取れたと
    言える。
 
    $$ H \Psi = \bigoplus (i \hbar ^{\nabla})^{\oplus L} $$
    $$ = \bigoplus {{ H \Psi} \over \nabla L} $$
    $$ = e^{x \log x} = x^{(x)^{'}} $$
    however,
    $$ {d \over df}F = m(x) $$
    and gravity equation in being abled to existed with quantum physics,
    and this physics is also represented with quantum level of differential
    geometry.
    $$ {d \over dt}\psi(t) = \hbar $$
    $$ = {1 \over 2}i e^{i \Hat{H}} $$
    $$ (i \hbar)^{'} = (- e^{i \Hat{H}})^{'} $$
    $$ = - i e^{i \Hat{H}} $$
    $$ \psi(x) = e^{-i \Hat{H}t}, \bigoplus (i \hbar ^{\nabla})^{\oplus L} 
    = {{{1 \over 2} e^{i \Hat{H}}}^{(-i e^{i \Hat{H}})}} $$
    $$ = ({1 \over 2}f)^{-i f}$$
    $$ = ({1 \over 2})^{-i f} \cdot e^{-x \log x} \cdot (f)^{i} $$
    $$ = \int e^{- x} x^{t - 1}dx, {d \over d \gamma}\Gamma = e^{-x \log x}  $$
     
    $$ f = x, i = t, {1 \over 2} = a, \bigoplus a^{-t x}x^t [I_m]
     \cong \int e^{-x} x^{t - 1}dx $$
     Quantum level of differential geomerty is also resulted with
    Gamma function of global differential manifold of values, Heisenberg
    equation is alos resulted with global integrate manifold and this
    quantum level of differential geometry is manifold of deprive of formula.
    $$ |\psi(t)\rangle_s = e^{-i\Hat{H} t}|\Psi\rangle_{H},
    \Hat{A}_s = \Hat{A}_H(0) $$
    $$ |\Psi(t)\rangle_s \to {d \over dt} $$
    $$ i {d \over dt}|\psi (t) \rangle_s = \Hat{H}|\psi(t)\rangle_s $$
    $$ \langle\Hat{A}(t)\rangle = \langle\Psi(t)|\Hat{A}(0)|\Psi(t)\rangle  $$
    $$ {d \over dt}\Hat{A} = {1 \over i}[\Hat{A},H] $$
    $$ \Hat{A}(t) = e^{i\Hat{H}t}\Hat{A}(0)e^{-i\Hat{H}t} $$
    $$ \lim_{\theta \to 0} \begin{pmatrix} \sin \theta \\
    \cos \theta \end{pmatrix} \begin{pmatrix} \theta & 1 \\
    1 & \theta \end{pmatrix} \begin{pmatrix} \cos \theta \\
    \sin \theta \end{pmatrix} = \begin{pmatrix} 1 & 0 \\
    0 & - 1 \end{pmatrix} $$
    $$ f^{-1}(x) x f(x) = I^{'}_m, I^{'}_m = [1,0] \times [0,1] $$

    $$ {x + y} \ge \sqrt{x y} $$
    $$ {x^{{1 \over 2} + iy} \over e^{x \log x}} = 1 $$
    $$ \mathcal{O}(x) = \nabla_i \nabla_j \int e^{{2 \over m}\sin \theta 
    \cos \theta } \times {{N \mathrm{mod}(e^{x \log x})} \over {O(x)(x +
    \Delta |f|^2)^{1 \over 2}}} $$
    $$ x \Gamma(x) = 2 \int|\sin 2 \theta|^2 d \theta, \mathcal{O}(x)
    = m(x)[D^2 \psi] $$
    $$ i^2 = (0,1) \cdot (0,1), |a||b|\cos \theta = - 1 $$
    $$ E = \mathrm{div}(E, E_1) $$
    $$ \left({\{f,g\}} \over [f,g] \right) = i^2, E = mc^2, I^{'} = i^2 $$
     Gamma function of global differential equation is first of differential
    function deprivate with ordinary differential equations, more also
    universe and other dimensiion of Higgs field emerge with zeta function,
    more spectrum focus is three dimension of manifold of singularity or pair
    theorem, and universe and other dimension of each value equals.
    This reverse of circle function of manifolds is also Higgs fields of 
    dense of entropy and result equation. And differential structure of quantum
    level in gravity is poled with zeta function and quantum group in high 
    result values. Zeta funtion and quantum group is Global differential
    manifold of result in values.
    $$ {d \over d \gamma} \Gamma = m(x) $$
    $$ = e^{- x \log x} $$
    $$ \sin i x = {{e^{-x} + e^{x}} \over 2i} $$
    $$ {d \over df}F = m(x) = e^{f} + e^{-f} $$
    $$ = 2 i \sin (i {x \log x}) $$


   Beta function is,
      $$ \beta(p,q) = \int x^{1 - t} (1 - x)^{t}dx = \int t^{x}(1 - t)^{x - 1}
      dt $$
      $$ 0 \le y \le 1, \int^{1}_{0} x^{10}(1 - x)^{20}dx = B(11,21) $$
      $$ = {{\Gamma(11)\Gamma(21)} \over \Gamma(32)} = {{10! 20!} \over {
	      31!}} = {1 \over {931395465}} $$
       $$ {1 \over {931395465}} = {1 \over 9} = {1 \over {1 - x}} $$
	 $$     = {1 \over {1 - z}} =
	      \sum_{k=0}^{\infty}z^k = {1 \over {1 + z^2}} = \sum_{
		      k=0}^{\infty} (-1)^k z^{2k} $$
      $$ f(x) = \sum_{k=0}^{\infty} a_k z^k $$
      $$ {{d^n y} \over dx^n} = n! y^{n+1} $$
      $$ f^{(0)}(0) = n!f(0)^{n+1} = n! $$
      $$ f(x) \cong \sum_{k=0}^{\infty}x^n = {1 \over {1 - x}} $$
      In example script is,
      $$ {{dy} \over dx} = y^2, {1 \over y^2}\cdot {dy \over dx} = 1 $$
      $$ \int {1 \over y^2}{dy \over dx}dx = \int{dy \over y^2} = - {1 \over y
      } $$
      $$ - {1 \over y} = x - C, y = {1 \over {C - x}} $$
      $$ \exists x = 0, y = 1 $$
      $$\text{in first value condition compute with} $$
      $$ \text{result, C consumer sartified,} $$
      $$ y = {1 \over {1 - x}} $$
      This value result is concluded with native function from Abel manifold.

      Abelian enterstane with native funtion meltrage from Daranvelcian
      equation on Beta function.
 

   $$ H\Psi = \bigoplus {H \Psi \over \nabla L} $$
  　
 $$ \Gamma = \int e^{-x}x^{1 - t}dx $$
   $$ \gamma = \int e^{-x}x^{1 - t}\log x dx $$
   $$ = \int \Gamma(\gamma)^{'} dx_m$$
   $$ = {d \over d \gamma}\Gamma $$

   Eight differential geometry are each intersect with own level of
   concept from expalanation of Euler product system.
   This component of three manfold sergeried with geometry of destroy and
   desect with time element of Stokes equation.
   ８種類の微分幾何では、それぞれの時間の固有値が違う。
   閉3次元多様体上で微分幾何の切除、分解によって時間の性質が決まる。
   This manifold gut theory from described with zeta function to
   catastrophe for non tree of routs result on sergery of space system.
   閉３次元多様体に統合されると、ゼータ関数になり、各幾何に分解された場合
   に、非分岐から、この上、sergeryの結果が決まる。

   各微分構造においての時間発展での熱エネルギーの変化は$E=mc^2, E=m^2 c^2 $
   のthree manifoldが微分幾何構造の変化にそれぞれ対応する。
   Each geometry point interacte with exchange of three manifold in
   Seifert structure from special relativity of equation for
   heat energy fluentations on time developyed from this extermaination.

   $$ \nabla ({i \hbar^{\nabla}})^{\oplus L}, \bigoplus ({i \hbar^{\nabla}})^{
	   \oplus L}, \square ({i \hbar^{\nabla}})^{\oplus L} $$
   These operator of equation on summatative manifolds from emerged with
   element of particle conclution, Euler product is resulted from this
   operator expalanation of locality insectations.
   これらの作用素が加減乗除の生成式の元、Euler Productの結論による
   作用素生成の論理素子でもある。
   Moreover, this eight geometry of differential operator are
   constructed with four pattern of Jones manifold from
   summatative formula and this system extate with special relativity
   references. This decieved of elemet on summatative equation routed
   to internext in real and imaginary pole on complex dimension,
p  and this dimension explation with Stokes theorem defined too.
   And this defined circutation of Yacobi matrix is represected with
   Knot theory with anstate with Jones manifold.
   その上に、8種類の微分幾何は、Jones多項式の4パターンでの構成される
   差分と加群方程式から、特殊相対性理論をも思わせる、この差分方程式
   $$ {d \over d \gamma}\Gamma = e^{f} + e^{-f} \ge e^{f} - e^{f} $$
   から、8種類の代数幾何が、ストークスの定理と同じく、このJones多項式
   の固有周期が、すべての時間の流れの基軸をも内包してもいることが、
   わかる。このJones多項式は、結び目の理論をも、この周期が言い表してもいる。
   $$ {d \over d \gamma}\Gamma + \int C dx_m = 2(\cos(i x \log x) - i \sin
   (i x \log x)) $$
    $$ e^{- \theta} = \cos(i x \log x) - i \sin(i x \log x)$$
    $$ \int \Gamma(\gamma)^{'}dx_m = 2(\cos (i x \log x) 
    - i \sin(i x \log x)) $$

 カルーツァ・クライン空間の方程式は、
    $$ ds^2 = g_{\mu\nu}(x)dx^{\mu}dx^{\nu} + \psi^2(x)(dx^2 + \kappa^2
    A_{\mu}(x)dx^{\nu})^2 $$
    この式は、
    $$ ds^2 = -N(r)^2dt^2 + \psi^2(x)(dr^2 + r^2d\theta^2) $$
    $$ ds^2 = -dt^2 + r^{-8\pi Gm}(dr^2 + r^2d\theta^2) $$
    とまとまり、
    $$ dx^2 = g_{\mu\nu}(x)(g_{\mu\nu}(x)dx^2 - dxg_{\mu\nu}(x)) $$
    $$ dx = (g_{\mu\nu}(x)^2dx^2 - g_{\mu\nu}(x)dx g_{\mu\nu}(x))^{1 \over 2} $$
    と反重力と正規部分群の経路和となり、　
    $$ \pi(\chi, x) = i \pi (\chi, x)f(x) - f(x)  \pi (\chi, x) $$
    と基本群にまとまる。
    シュワルツシルト半径は、
    $$ ds^2 = -(1 - {r_s \over r})c^2 dt^2 + {1 \over {{1 - {r_s \over r}}}} 
    dr^2 + r^2 d \theta^2 + r^2 \sin^2 \theta d \psi^2 $$
    これは、カルーツァ・クライン空間と同型となる。
    一般相対性理論は、
    $$ R^{\mu\nu} + {1 \over 2}g_{ij}\Lambda = \kappa T^{\mu\nu} $$
    この多様体積分は、
    $$ \int \kappa T^{\mu\nu}\mathrm{dvol} = \int  \left(R^{\mu\nu} + 
    {1 \over 2}g_{ij}\Lambda\right)\mathrm{dvol}  $$
    ガンマ関数のおける大域的微分多様体は、これと同型により、
    $$ \int \Gamma \cdot {d \over {d \gamma}}{\Gamma dx_m} \le \int \Gamma
    dx_m + {d \over d \gamma}\Gamma $$
    $$ = \int \Gamma(\gamma)^{'}dx_m $$
    オイラーの定数の多様体積分は、
    $$ \int \left({{\int {1 \over x^s}}dx - \log x} \right)\mathrm{dvol} $$
    この解は、シュワルツシルト半径と同型より、
    $$ = e^f - e^{-f} \le {e^f + e^{-f}} $$
    次元の単位は多様体より、大域的微分多様体とオイラーの定数
    の多様体積分の加群分解は、オイラーの公式の三角関数の虚数の度解より、
    $$ {d \over df}F + \int C dx_m = 2 (\cos (i x \log x) 
    - i \sin (i x \log x)) $$
    すべては、大域的微分多様体の重力と反重力方程式に行き着き、
    $$ {d \over df}F \ge {d \over df}\int \int{1 \over (x \log x)^2}dx_m
    + {d \over df}\int \int{1 \over (y \log y)^{1 \over 2}}dy_m $$
    と求まる。


  ８種類の微分幾何での、それぞれの時間の固有値の動静は、
  ストークスの流体力学での、熱エネルギーの各流れの仕方で表されている。
  種数１の2種が種数０の球体に両辺で交わっていると、
一般相対性理論における重力理論は大域的微分多様体と積分多様体についての 単体量を求めるための空間の対数関数における不変性を記述するプランクスケール と異次元の宇宙におけるスケール、ウィークスケールと言われているanti-D-brane として、この２種類の計量をガンマ関数とベータ関数としてオイラーの定数 とそれによる連分数が微分幾何の量子化と数式の値を求めると同じという 予知と推測値から、広中平祐定理の４重帰納法のオイラーの公式による 多様体積分と同類値として、ペレルマン多様体がサーストン空間に成り立つ と同じく、広中平祐定理がヒルベルト空間で成り立つとしての、 これら２種の多様体の場の理論が、ゼータ関数がAdS5多様体で成り立つことと、 不変量として、ゼータ関数がこの3種の多様体の場の理論のバランスをとる
  理論として言えることである。
$$ ||ds^2|| = \lim_{x \to \infty} [\delta(x) \int \int \int \pi \left( \sum_{k=0}^{\infty}
		{{}^n\sqrt{p},x \over n} \right)
^{1 \over 2}d \tau]^{\mu\nu} $$
  This norm is component of fields in Hilbert manifold of space theorem rebuilt with
Mebius space in Gamma Function on Beta function fill into power of rout fundamental
group.

  　 Gravity of general relativity theory describe with Cutting of space
in being discatastrophed from Global differential and integral manifold
of scaled levetivity in plank scal and the other vector of universe
scale, these two scale inspectivity of Gamma and Beta function escourted
with these manifold experted in result of Differentail geometry of
quantum level manifold equal with Euler product and continue parameter,
moreover this product is relativity of Hironaka theorem in Four assembled
integral of Euler equation.

 重力場理論の式は、
  Gravity equation is
  $$ \square = - {16 \pi G \over c^4}T^{\mu\nu} $$
  This equation quated with being logment of formula, and
  this formula divided with universe of number in prime zone,
  therefore, this dicided with varint equation is monotonicity of
  being composited with Weil's theorem united for Gamma function.
  この式は、対数を宇宙における数により求める素数分布論として、この
  大域的積分分断多様体がガンマ関数をヴェイユ予想を根幹とする単体量
  として決まることに起因する商代数として導かれる。
  $$ \scalebox{1}[2]{\alearletter} = 
  {{{{8 \pi G} \over {c^4}}T^{\mu\nu}}/ \log x} $$
  $$ {}^{t}\scalebox{1}[2]{\variiint}\mathrm{cohom D_{\chi}[I_m]}$$
  $$ = \oint {(px^n + qx + r)^{\nabla l}} $$
  $$ {d \over d l}L(x,y) = 2 \int ||\sin 2x||^{2}d \tau $$
  $$ {d \over d \gamma}\Gamma $$
  この関数は大域的微分多様体としてのアカシックレコードの合流地点として、
  タプルスペースを形成している。
  This function esterminate with acasic record of global differential
  manifolds.
  $$ = [i \pi (\chi, x), f(x)] $$
  それにより、この多様体は基本群をアカシックレコードの相対性としての
  存在論の実存主義から統合される多様体自身としてのタプルスペースの
  池になっている。
  And this manifold from fundemental group esterminate with
  also this manifold estimate relativity of acasic record.
  $$ {d \over d \gamma}\Gamma = \int C dx_m = \int(\int {1 \over x^s}dx
  - \log x)\mathrm{dvol} $$
  また、このアカシックレコードはオイラーの定数のラムダドライバー
  にもなっている。
  More also this record tupled with lake of Euler product.

\title{特殊相対性理論の虚数回転による多様体積分と、\\
それによる一般相対性理論の再構築理論}
\author{}
   
  \newcommand{\oiint}{%
	   \begin{picture}(15,15)
		   \put(0,0){
			   $ \scalebox{1}[2]{$\iint$} $}
		   \put(2,2){$\bigcirc$}
	   \end{picture} %
	   }


 $$ {\partial \over \partial f}F(x) = \int \int \text{cohom}D_k(x)[I_m]  $$
 Cohomology element is constructed with isotopy of D-brane,
 and this isotopy of symmetry structure is sheaf of manifold,
 more over this brane of element is double integrate of projection.
 $$ {\partial \over \partial f}F = {{}^t \scalebox{1}[2]{\variint}}
 \text{cohom}D_k(x)^{\ll p} $$

  $$  {\partial \over \partial f}F = {{}^t \scalebox{1}[2]{\variint}}
 \text{cohom}D_k(x)^{\ll p} = 
\scalebox{1}[2]{\alearletter}$$ 
   $$ \nabla_{i} \nabla_{j} \int f(x) d \eta = {\partial^2 \over {\partial x
   \partial y}}\int \scalebox{1}[2]{\alearletter} d \eta $$
   一般相対性理論の加群分解が偏微分方程式と同じく、
   特殊相対性理論の多様体積分の虚数回転体がベータ関数となる。
   ほとんどの回転体の体積が、係数と冪乗での回転体として、ベータ関数
   と言える。
   $$ = R^{{\mu\nu}^{'}} + {1 \over 2}\Lambda g_{ij}^{'} = \int 
   \left(i {{v \over {
	   \sqrt{1 - 
   {{{({v \over t})}}^{2}}}}}} + {v \over {\sqrt{1 - 
   {{{({v \over t})}}^{2}}}}}\right)\mathrm{dvol} $$
   $$ = {d \over df}\int \int{1 \over {(x \log x)^2}}dx_m +
   {d \over df}\int \int {1 \over {(y \log y)^{1 \over 2}}}dy_m $$ 
    この大域的積分多様体が大域的微分多様体の反重力と重力方程式
    で表せられて、
   $$ = \int C dx_m = \int \kappa T^{\mu\nu} \mathrm dx_m = T^{{\mu\nu}{^{T^
   {\mu\nu}}}^{'}}$$
   オイラーの定数の大域的積分多様体が、一般相対性理論の大域的積分多様体
   であり、エントロピー不変式で表せられる。

  $$ \scalebox{1}[2]{\alearletter} = 
  {{{{8 \pi G} \over {c^4}}T^{\mu\nu}}/ \log x} $$
  $$ {}^{t}\scalebox{1}[2]{\variiint}\mathrm{cohom D_{\chi}[I_m]}$$
  $$ = \oint {(px^n + qx + r)^{\nabla l}} $$
  $$ {d \over d l}L(x,y) = 2 \int ||\sin 2x||^{2}d \tau $$
  $$ {d \over d \gamma}\Gamma $$
 $$ ||ds^2|| = \lim_{x \to \infty} [\delta(x) \int \int \int \pi \left( \sum_{k=0}^{\infty}
		{{}^n\sqrt{p},x \over n} \right)
^{1 \over 2}d \tau]^{\mu\nu} $$

 $$\oiint  \cong  ||ds^2|| = \lim_{x \to \infty} 
 [\delta(x) \int \int \int \pi \left( \sum_{k=0}^{\infty}
		{{}^n\sqrt{p},x \over n} 
		\right)^{1 \over 2}d \tau]^{\mu\nu} $$
 
 $$ e^{-2\pi T||\psi||}[\eta + \bar{h}]dx^{\mu}dx^{\nu} + T^2 d^2 \psi $$
 $$ {{-16 \pi G} \over c^4}T^{\mu\nu}/ \log x =  
 \scalebox{1}[2]{\alearletter} $$
 $${{-16 \pi G} \over c^4}T^{\mu\nu}/ e^{-2\pi T||\psi||} $$
 $$ = 4 \pi G \rho $$
 $$ {\partial \over {\partial x \partial y}} = \nabla_{i} \nabla_{j} $$ 
 $$ \square \iiint  = {}^{t}\scalebox{1}[2]{\variiint}\ $$
 $$ {\partial \over {\partial x}} \iiint = \nabla_{i}\nabla_{j} \int
 \nabla f(x) d \eta $$
 $$ {\partial \over {\partial f}}\iiint = \square \iiint $$

 $$ \reflectbox{$\Gamma$} | \Gamma, \reflectbox{B} | B $$
 $$ \reflectbox{E} | E, \reflectbox{C} | C $$
 $$ \reflectbox{F} | F, \reflectbox{$\beta$} | \beta $$
 $$ \reflectbox{D} | D, {}^{t}\scalebox{1}[2]{\variint} \cong 
 \bigoplus D^{\bigoplus L} $$
 $$ \square +  \scalebox{1}[2]{\alearletter} = \emptyset $$
 $$  \scalebox{1}[2]{\alearletter} | \emptyset = \square $$
 $$ {\begin{pmatrix} 1 & 0 \\
 0 & -i \end{pmatrix}}^{1 \over 2} \Biggl| {\begin{pmatrix} 1 & 0 \\
 0 & -i \end{pmatrix}}$$
 $$ x^{1 \over 2} \cong x, \emptyset^{1 \over 2} = \emptyset$$
 
    自分がいると、方位が決まる。そうしたら、空間ができる。北と南が出来る。
    これは、磁気単極子が相対空間での根本的なアイデアらしい。
    無になると関係性がなくなり、磁気単極子もある上に、全て無でもある。
    無になると十二支縁起になる。宇宙と異次元の外が無でもある。
    運動量が決まると位置が決まる。位置だけだと宇宙が不確定性原理
    であり、異次元が合わさると運動量も位置も決まる。
    無になると位置も運動量も同時にわかる。
    このアイデアが3次元多様体のエントロピー値である。
    
    $$ \int^{\ll n} ||ds^2|| = \int\int\int 8 \pi G \left({p \over c^3} + 
    {V \over S}\right)dxdydz $$
    $$ = {\partial \over \partial f_n}{||ds^2||} $$


 $$ {\partial \over \partial f}F(x) = \int \int \text{cohom}D_k(x)[I_m]  $$
 Cohomology element is constructed with isotopy of D-brane,
 and this isotopy of symmetry structure is sheaf of manifold,
 more over this brane of element is double integrate of projection.
 $$ {\partial \over \partial f}F = {{}^t \scalebox{1}[2]{\variint}}
 \text{cohom}D_k(x)^{\ll p} $$
 And global partial defferential equation is integrate of cut in cohomolgy,
 and this step of defferential D-brane of sheap value is varint of equals,
 inverse of horizen of cute of section value is reverse of integrate of 
 operators. Three dimension of varint of manifolds in operator of sheaps,
 and this step of times of value is equal with D-brane of equations.
 Global differential equation and integrate equation is operator of 
 global scales of horizen cut and add of cut equation are routs.

 $$ \oint r dx_m = \mathcal{O}(x,y) $$

 $$ {{}^t \scalebox{1}[2]{\variiint}}_{{D(\chi,x)}} 
 \text{Hom}[D^2 \psi]^{\ll p} \cong 
 \text{vol}\left({V \over S}\right) $$

 $$ {\partial \over \partial f}F(x) = {{}^t \scalebox{1}[2]{\variint}} 
 \text{cohom}D_k(x)^{\ll p} $$

 $$ {d \over df}F = (F)^{f^{'}}, \int F dx_m = (F)^{f}  $$
 $$ {d \over df}\int \int{1 \over (x \log x)^2}dx_m =
 \left(\int \int{1 \over (x \log x)^2}dx_m \right)^{f^{'}}  $$
 $$ f^{'} = (2x (\log x + \log (x + 1))) $$
  $$  \int \int F dx_m = \left({1 \over (x \log x)^2}\right)^{
	  (\int 2x (\log x + \log (x + 1))dx)} $$ 
   $$= e^{-f} $$
  $$ {d \over df}F = \left({1 \over (x \log x)^2}\right)^{
	  (2x (\log x + \log (x + 1))) } $$  
  $$= e^{f} $$
  $$ \log (x \log x) \ge 2(y \log y)^{1 \over 2} $$

  $$ \log (x \log x) \ge 2(\sqrt{y \log y}) $$

 
    $$ ||ds^2|| = 8 \pi G \left({p \over c^3} + {V \over S}\right) $$
    $$ \int ||ds^2||dx_m  = \int 8 \pi G \left({p \over c^3} 
    + {V \over S}\right)\mathrm{dvol}  $$
    $$ \int ||ds^2||dx_m  = \int 8 \pi G \left({p \over c^3} 
    + {V \over S}\right)dy_m $$
     $$ \int ||ds^2||dx_m  = \int {1 \over {8 \pi G \left({p \over c^3} 
    + {V \over S}\right)}}dx_m $$
    $$ {d \over df}F = m(x), \bigoplus \left({i \hbar}^{\nabla}\right)^{
	    \oplus L} $$
    $$ = \nabla_{i} \nabla_{j} \int \nabla f(x)d \eta $$

    $$ dx_m = {y \over {\log x}}, dy_m = {x \over \log x} $$
    $$ e^{-f}dV = dy_m = \mathrm{dvol} $$
    偏微分は加群分解と同じ計算式に行き着く。

   宇宙と異次元の誤差関数のエネルギー、
    $AdS_5$多様体がベータ関数となる値の列が、異次元への扉となっている。
    $$ \beta(p,q) = \text{誤差関数} + \text{Abel多様体} $$
    $$ = \text{$AdS_5$多様体} $$
    $$ = {d \over df}F + \int C dx_m = \int \Gamma(\gamma)^{'}dx_m $$
    ここで、アーベル多様体はEuler productである。
    ベータ関数の数列がわかると、ゼータ関数は無であるというのが、
    どういうことかが、物、物体に影ができて、ものが瞑想と同じであり、
    これから、風景がベータ関数の数値列に見えるらしい。
     この大域的微分多様体のガンマ関数が、複素力学系のマンデルブロ集合の
    プリズムと同じ構造の見方らしい。

    $$ ||ds^2|| = e^{-2\pi G ||\psi||}[\eta + \bar{h}(x)]dx^{\mu}dx^{\nu}
    + T^2 d^2 \psi $$
    $$ \beta(p,q) = \text{誤差関数} + \text{Abel多様体} $$
    $$ \int \mathrm{dvol} = \square \psi $$
    $$ \int \nabla \psi^2 d \nabla \psi = \square \psi $$
    $$ \text{expanding of universe} = \text{exist of value} $$
    $$ = \log (x \log x) = \square \psi $$
    $$ \text{freeze out of universe} = \text{reality of value} $$
    $$ = (y \log y)^{1 \over 2} = \nabla \psi $$

    All of value is constance of entropy, universe is freeze out
    constant, and other dimension is expanding into fifth dimension of
    inner.


      $$ x^n + y^n = z^n, \beta(p,q) = x^n + y^n - \delta(x) = z^n - \delta
      (x) $$
      $$ {d \over dt}g_{ij} = - 2 R_{ij}, {\cos x \over (\cos x)^{'}}
      \cdot (\sin x)^{'} = z_n, z^n = - 2 e^{x \log x} $$

      $$ z = e^{-f} + e^{f} - y $$
      $$ \beta(p,q) = e^{-f} + e^{f} $$
      相対性は、暗号解読と同じ仕組みの数式を表している。
      ここで言うと、yが暗号値である。チェックディジットと同じ仕組みを有
      している。

      $$ \square x = \int {f(x) \over {\nabla(R^{+}\cap E^{+})}}d \square x$$
      $$= \int {{\Delta f(x) \circ E^{+}} \over {\nabla(R^{+}\cap E^{+})}}
      \square x $$
      $$ \square x = \int {d(R \nabla E^{+}) \over {\nabla(R^{+} \cap E^{+})}}
      d \square x $$
      $$ = \int {{\nabla_{i} \nabla_{j} (R + E^{+}} \over {\nabla(R^{+}
      \cap  E^{+})}}d \square x $$
      $$ x^n + y^n = z^n $$
      $$ \mathrm{exp}(\nabla(R^{+}\cap E^{+}),\Delta(C \supset R)) $$
      $$ = \pi (R_1 \subset \nabla E^{+}) = \mathrm{rot}(E_1, 
      \mathrm{div}E_2) $$
      $$ xf(x) = F(x), s\Gamma(s) = \Gamma(s+1) $$
      $$ Q \nabla C^{+} = {d \over df}F(x) \nabla \int \delta(s)f(x)dx $$
      $$ E^{+}\nabla f = {{e^{x \log x} \nabla n! f(x)} \over {E(x)}} $$
      $$ {d \over df}F ~ F^{{f}^{'}} = e^{x \log x} $$
      $$ (C^{\nabla})^{\oplus Q} = e^{x \log x} $$
      $$ {d \over dt}g_{ij} = - 2 R_{ij}, {\cos x \over (\cos x)^{'}}
      \cdot (\sin x)^{'} = z_n, z^n = - 2 e^{x \log x} $$

      $$ R \nabla E^{+} = f(x) \nabla e^{x \log x}, {d \over df}F = F^{
	      {f}^{'}} = e^{x \log x}  $$
      南半球と単体（実数）の共通集合の偏微分した変数をどのようなF(x)かを
      $$ \int \delta(x)f(x)dx $$
      と同じく、単体積分した積分、共通集合の偏微分をどのくらいの微分変数
      を
      $$ \int \nabla \psi^2 d \nabla \psi = \square \psi $$
      と同じ、
      $$ \int dx = x + \text{C(Cは積分定数) }$$
      と原理は同じである。


      Beta function is,
      $$ \beta(p,q) = \int x^{1 - t} (1 - x)^{t}dx = \int t^{x}(1 - t)^{x - 1}
      dt $$
      $$ 0 \le y \le 1, \int^{1}_{0} x^{10}(1 - x)^{20}dx = B(11,21) $$
      $$ = {{\Gamma(11)\Gamma(21)} \over \Gamma(32)} = {{10! 20!} \over {
	      31!}} = {1 \over {931395465}} $$
       $$ {1 \over {931395465}} = {1 \over 9} = {1 \over {1 - x}} $$
	 $$     = {1 \over {1 - z}} =
	      \sum_{k=0}^{\infty}z^k = {1 \over {1 + z^2}} = \sum_{
		      k=0}^{\infty} (-1)^k z^{2k} $$
      $$ f(x) = \sum_{k=0}^{\infty} a_k z^k $$
      $$ {{d^n y} \over dx^n} = n! y^{n+1} $$
      $$ f^{(0)}(0) = n!f(0)^{n+1} = n! $$
      $$ f(x) \cong \sum_{k=0}^{\infty}x^n = {1 \over {1 - x}} $$
      In example script is,
      $$ {{dy} \over dx} = y^2, {1 \over y^2}\cdot {dy \over dx} = 1 $$
      $$ \int {1 \over y^2}{dy \over dx}dx = \int{dy \over y^2} = - {1 \over y
      } $$
      $$ - {1 \over y} = x - C, y = {1 \over {C - x}} $$
      $$ \exists x = 0, y = 1 $$
      $$\text{in first value condition compute with} $$
      $$ \text{result, C consumer sartified,} $$
      $$ y = {1 \over {1 - x}} $$
      This value result is concluded with native function from Abel manifold.

      Abelian enterstane with native funtion meltrage from Daranvelcian
      equation on Beta function.
 
       Rotate with pole of dimension is converted from 
       other sequence of element,this atmosphere of vacume 
       from being emerged with quarks of being constructed
with Higgs field, and this created from energy of zero dimension are begun with
universe and other dimension started from darkmatter that represented with 
big-ban.
Therefore, this element of circumstance have with quarks of twelve lake with
atoms.

  $$ (dx, \partial x)\cdot (\epsilon x, \delta x) = \begin{pmatrix}
	  1 & 0 \\
  0 & - 1 \end{pmatrix}^{1 \over 2}  $$
  $$ {d \over df} F = m(x) $$

 And this phenomounen is super symmetry theorem built with quarks of element, 
also this chemistry of mechanism estimate from physics of operatsion.
These response of mechanism chain of geometry into space being emerged with
creature and univse of existing of combination.

  This converted with dimension of element also emerged from imaginary and 
reality of pole's space.

  And this pole of transport of dimension belong with vector of time has 
with one rout of sequecence. Quanum physics also belong with other vector 
of time has with imaginary rout of sequecense.

 This sequence of being estimated with non fluer of time, and this space 
of element have with gravity and antigravity of power.
Other vecor of time is antigravity rout of sequecense.

 Weak electric theory is estimate from time has with one rout of sequecense,
this topology of chain is resulted from time of rout ways.

      $$ \square ({{\sigma_1 + \sigma_2} \over 2 }) = [3 \pi (\chi,x) \circ f(x),
   \sigma(x)] $$
   $$ \square \psi = 8 \pi G T^{\mu\nu} $$
   $$ {d \over dt}g_{ij}(t) = - 2 R_{ij} $$
   $$ \nabla \psi^2 = 4 \pi G \rho $$
   $$ \square (\sigma_1 + \sigma_2) = [6\pi(\chi,x)\circ f(x), \sigma(x)] $$
   $$ = [i \pi(\chi,x),f(x)] $$
   $$ \square \psi = [12 \pi (\chi,x)\circ f(x), \sigma(x)] $$
   $$ {d \over df} F = \int e^{-f}[- \Delta v + R_{ij} v_{ij}+ 
   \nabla_i \nabla_j f  + v \nabla_i \nabla_j + 2<f,h> + (R + \nabla f)
   ({v \over 2} - h) ] $$
   $$ =[i \pi(\chi,x), f(x)] $$

   These equation is reminded time pass rout of one rout way of forms,
   and this rout of time ways which go for system from future and past.
   Therefore, this resulted system of time mechanism is one true flow
   that weak electric theorem oneselves.
   and moreover, one rout time way of forms is reverse with antigravity of
   time system. This also spectrum focus is true that Maxwell theorem and
   strong boson unite with antigravity, this unite is essense on the contrary
   from weak electric theorem, this theorem called for time rout forms is
   strong electric theorem.
   This two theorem is united with quantum physics that no time flow system.

   $$ f^{-1}(x)xf(x) = 1,H_m = E_m \times K_m $$


   The non-commutative theorem is constructed from world line surface that
   this complex manifold estimate with rolanz atractor, and this string
   theorem have with one world of universe mate six quarks and other world of
   dimension mate other element of six quarks.
   These particle quarks is built with super symmetry space of dimension.

   $$ i = (1,0)\cdot (1,0), e^{i\theta} = \cos \theta + i \sin \theta $$
   $$ e^{-i \theta} = \cos \theta - i\sin \theta $$
   $$ \sin \theta = {e^{i \theta} - e^{-i \theta} \over 2i} $$
   $$ \sin i \theta = {e^{-\theta} - e^{\theta} \over 2} $$
   $$ \pi (\chi,x) = \cos \theta + i \sin \theta $$

   In this equations, two dimension redestructed into three dimension,
   this destroy of reconstructed way is append with fifth dimension.
   This deconstructed way of redestructe is arround of universe attached
   with three dimension, this over cover call into fifth dimension.

   $$ R(-\alpha) = \begin{pmatrix} \cos \alpha & \sin \alpha \\
   - \sin \alpha & \cos \alpha \end{pmatrix} $$
   $$ R(\alpha) = \begin{pmatrix} \cos \alpha & - \sin \alpha \\
   \sin \alpha & \cos \alpha \end{pmatrix} $$
   $$ R(\alpha)M R(-\alpha) = 
   \begin{pmatrix} \cos \alpha & - \sin \alpha \\
   \sin \alpha & \cos \alpha \end{pmatrix} \begin{pmatrix} 1 & 0 \\
   0 & - 1 \end{pmatrix} \begin{pmatrix} \cos \alpha & \sin \alpha \\
   - \sin \alpha & \cos \alpha \end{pmatrix} $$
   $$ = \begin{pmatrix} \cos^2 \alpha - \sin^2 \alpha & 2 \sin \alpha \cos 
	   \alpha \\
	   2 \sin \alpha \cos \alpha & - \cos^2 \alpha + \sin^2 \alpha
   \end{pmatrix} $$
   $$ = \begin{pmatrix} \cos 2 \alpha & \sin 2 \alpha \\
   \sin 2 \alpha & - \cos 2 \alpha \end{pmatrix} $$

   $$ \begin{pmatrix} \cos \alpha & - 1 \\
	   1 & - \sin \alpha \end{pmatrix} \begin{pmatrix} \cos \alpha \\ 
		   \sin \alpha
   \end{pmatrix} = \begin{pmatrix} 1 & 0 \\
   0 & - 1 \end{pmatrix} $$
   $$ \begin{pmatrix} \cos \alpha & \sin \alpha \\
   \sin \alpha & - \cos \alpha \end{pmatrix} \begin{pmatrix}
   x \\ y \end{pmatrix} = \begin{pmatrix} 1 & 0 \\
   0 & - 1 \end{pmatrix} $$
   $$ \lim_{\theta \to 0} \begin{pmatrix} \sin \theta \\ \cos \theta 
   \end{pmatrix} \begin{pmatrix} \theta & 1 \\
   1 & \theta \end{pmatrix} \begin{pmatrix} \cos \theta \\ \sin \theta
	   \end{pmatrix} = \begin{pmatrix} 1 & 0 \\ 0 & - 1 \end{pmatrix} $$
		   $$ f^{-1}xf(x) = 1 $$
		   $$ (\log x)^{'} = {1 \over x}, x^n + y^n = z^n  $$
		   $$ x^n = - y^n + c, nx^{n-1} = -n y^{n-1}y^{'} $$
		   $$ y^{'} = {{n x^{n-1}} \over {n y^{n-1}}} $$
		   $$ = {x^{n-1} \over y^{n-1}}= - {y \over x}\cdot
		   ({x \over y})^n $$
		   $$ - {\cos x \over (\cos x)^{'}}(\sin x)^{'} = z_n $$
		   $$ z^n = - 2 e^{x \log x} $$
		   $$ \lim_{x \to \infty} f(x) = a, \lim_{y \to \infty}
		   f(y) = b, \lim_{x,y \to \infty}\{f(x) + f(y)\} = a + b  $$
		   $$ \lim_{z \to 0}f(z) = c, \delta \int z^n = {d \over dV}
		   x^3 $$
		   $$ \lim_{x \to \infty} = c - \lim_{y \to \infty}f(y)
		   \lim_{x \to \infty} f(x) + \lim_{y \to \infty}f(y) = 
		   \lim_{z \to \infty}f(z) $$
		   $$ z^n = \cos n \theta + i \sin n \theta $$
		   $$ = - 2 e^{x \log x} $$

  These equation is gravity and antigravity equation. 

	   $$ {d \over d\sigma}\begin{bmatrix}{(\sigma_1 + \sigma_2) 
	   \over 2}\end{bmatrix}  $$
	   $$ = \sigma(\upharpoonleft \downharpoonleft) + \sigma(
	   \upuparrows)+ \sigma(\downdownarrows)+ 
	   \sigma(\leftleftarrows)
	   + \sigma(\rightrightarrows) $$
	   $$ {d \over df}{\int \int {1 \over (x \log x)^2}}dx_m
	   = \sigma(\upharpoonright) $$
	   $$ {d \over df}{\int \int {1 \over (y \log y)^{1 \over2}}}dy_m
	   = \sigma(\downharpoonright) $$
	   $$ \sigma(\leftarrowtail) + \sigma(\rightarrowtail)
	   = \int e^{-f}[-\Delta v + R_{ij}v_{ij} + \nabla_{i}\nabla_{j}v
	   + v\nabla_{i}\nabla_{j} + 2 <f,h> + (R + \nabla f)(v - 
	   {h \over 2})]
	   $$
	   $$ \sigma(\downharpoonleft) = \sigma(\upharpoonright 
	   \downharpoonright + \upuparrows 
	   + \rightrightarrows)$$
	   $$ \sigma(\upharpoonleft) =  \sigma(\upharpoonright 
	   \downharpoonright + \downdownarrows 
	   + \leftleftarrows) $$
         
	   $$ \text{weak electric theorem} = \sigma(\upharpoonleft) $$
	   $$ \text{strong electric theorem} = \sigma(\downharpoonleft) $$

	   These equation is represented with topology of string model,
	   and weak electric theorem is constructed with Maxwell theorem 
	   and weak boson, gravity that estimate with this three power united.
	   Moreover, strong boson and Maxwell theorem, antigravity that
	   also estimate with this three power united.
	   This two united power is integrated from gravity and antigravity.
	   Then this united power is zeta function.

  Euler equation represet form,
  $$ C = \int {1 \over x^s}dx - \log x $$
  $$ = \int \left(\int {1 \over x^s}dx - \log x \right)\mathrm{dvol} $$
  $$ = \int \int{1 \over x^s}\mathrm{dvol} - \int  \log x \mathrm{dvol} $$
  This equation componet of $ \text{Higgs equation + Euler equation =
  Zeta function} $. This represent expression is Rich flow equation.
  This equation assent from Differential geometry in quantum level,
  Gamma and Beta function, Gauss surface, Higgs field, Einsetin tensor,
  Euler constance and so on.
    This equation is add and average of possibility, and universe, atom
  more over a certain theory, this universe oneselves is environment
  seek, other distance is possibility equation, however, other dimension
  cover demand with not only norm of non liner in three manifold
  but also non possibility in complete of environment equation.
    $$ = {d \over df} \int \int{1 \over (y \log y)^{1 \over 2}}dy_m 
  - \int e^{-x}x^{1 - t} \log x dx_m $$
  $$ \int x^{1 - t}e^{-x}dV = \int x^{1 - t}dm, \int x^{1 - t}e^{-x}dV
  = \int x^{1 - t}\mathrm{dvol}, f = \gamma = \Gamma^{'} = \int 
  e^{-x}x^{1 - t} \log x dx $$
  $$ = {d \over d \gamma}\Gamma^{-1} - (\gamma)^{\gamma^{'}} $$
  $$ = e^{-f} - e^{f} $$
  $$ = {d \over df}\int \int{1 \over (y \log y)^{1 \over 2}}dy_m
  - {d \over df}\int\int {1 \over (x \log x)^2}dx_m $$
  $$ = 2 \cos (i x \log x) $$
  $$  {d \over df}\int\int{1 \over (y \log y)^{1 \over 2}}dy_m
  + {d \over df}\int\int{1 \over (x \log x)^2}dx_m $$
  $$ = {d \over df}F = 2 i \sin (i x \log x) $$
  $$ {d \over df}F + \int C dx_m  
  = 2 (\cos( i x \log x) + i \sin(i x \log x)) $$
  $$ = 2 e^{- \theta} = 2 e^{- i x \log x} $$
  $$ = {d \over dt}g_{ij}(t) = - 2 R_{ij} $$
  $$ \log (x \log x) \ge 2 \sqrt{y \log y} $$
  $$ \log(x \log x) = \log x + \log \log x $$
  This universe and other dimension on blackhole and whitehole
  compont with universe and atom in movement and point of 
  environment value by a certain theory of complete environment 
  equation. And this three manifold of possibility equation 
  are able to research in both value formula.
  This envirnment of researchable reason is that surround of universe
  is setting to value by already in consented to being researched,
  then both value researchabled.
  And this three manifold equation is insert with Kaluza-Klein theory
  and Heisenbelg equation.
    $$ \nabla \psi^2 = 8 \pi G\left({p \over c^3} + {V \over S}\right) $$
  $$ \nabla \psi^2 = 8 \pi G \hbar + 8 \pi {V \over S} $$
  This energe of entropy value consist with whitehole of atom
  and universe in blackhole value, and this universe in balckhole 
  represent with Higgs filed, atom in whitehole consist with plack scale.
    $$ \square_v = 2 \sqrt{2 \pi G \hbar} + 2 \sqrt{2 \pi {V \over S}} $$
  This entropy value also consist with atom in balckole.
    $$ \sqrt{2 \pi T} = 2 \sqrt{2 \pi{V \over S}} $$

  These equation more also consist with universe being non gravity formula.
    $$ 2 \cos (i x \log x) + 2 i \sin(i x \log x) = 2 e^{-f} $$
  $$ 2 \sqrt{2 \pi {V \over S}} = {d \over df} 
  \int \int{1 \over (x \log x)}dx_m$$
  And also these equation consist with universe and other dimension in
  warped passeage in rout of integral.

   $$ e^{-2\pi ||psi|}[\eta_{\mu\nu} + \hbar{x}]dx^{\mu}dx^{\nu} 
  +  T^2 d^2 \psi $$
  $$ = e^{-f} + e^f = (u + d) + c, (w + s) + b = -e^{-f} + e^f $$
  この式は、時空にどれだけの原子が凝縮しているかを表していて、
  その上に、この逆数は、密度エネルギーをも示している。
  空間の体積にに原子を商代数にすると、濃度にもなる。
   逆数は、$\text{a=原子のエネルギー}$が全空間を1とすると
  $\text{$1 \over {y \over 
  x}$}$すると、$x \over y $は密度になる。　
  $\text{{x=原子},y=全空間、個数は}$$\text{$y \over x$}$ 
  $$ \square = \text{原子} \to {{[\eta_{\mu\nu} + \hbar(x)]dx^{\mu}dx^{\nu}
  } \over {e^{2 \pi T ||\psi||} }}$$
  $$ ||ds^2|| = e^{2 \pi T||\psi||}\left([\eta_{\mu\nu} + \hbar(x)]d^{\mu}dx^
  {\nu}\right)^{-1} + \left(T^2 d^2 \psi \right)^{-1} $$
  $$||ds^2|| = e^{-2 \pi T||\psi||}[\eta_{\mu\nu} + \hbar(x)]d^{\mu}dx^
  {\nu}] + T^2 d^2 \psi  $$
  $$ (u + d) + c = e^{-f} + e^f, (w + s) + b = e^{-f} + e^{f} $$
  $$ R_{ij}^{'} = - R_{ij} $$

  この$ AdS_5 $多様体の $e^{-2 \pi T ||\psi||} $は原子間距離を表していて、
  $ [\eta + \bar{h}(x)]d^{\mu}dx^{\nu} $で宇宙のD-braneの構造を表している。
  この数式が$T^2 d^2 \psi $とアーベル多様体が包み込んでいる。
  原子が回転するのと、ニュートンリングが回転する宇宙は同じ速さで回転する。　

  $$||ds^2|| = \square + \rho , ||ds^2|| = \nabla \Psi^2 + \psi $$
  $$ \eta_{\mu\nu} + \bar{h}(x) $$
  $$ \text{ploximetry} = \Delta x \Delta p - \Delta p \Delta x + 
  \delta(p,x)  $$
  $$ ||ds^2|| = e^{-2 \pi T ||\psi||} + [\eta_{\mu\nu}  + \bar{h}(x)]
  dx^{\mu\nu} + T^2 d^2 \psi $$
  $$ {d \over df}F = e^f + e^{-f}, {d \over df}\int C dx_m = e^f - e^{-f} $$
  $$ R_{ij}^{'} = - R_{ij} $$
  宇宙と異次元では０だが宇宙だけだと誤差が生じる。これが不確定性原理であり、
  異次元が誤差になり、宇宙と加群すると０になる。
  　$ \delta(p,x) $が隠れた変数となり、異次元で誤差をこの不確定性原理
  と表している。
  全ては素粒子方程式から生成される式達である。
   $$ (u + d) + c = e^{-f} - e^{f}, b + (w + s) = e^{-f} + e^{f} $$
   $$ \int \int{1 \over x^s}\mathrm{dvol} - \int \log x \mathrm{dvol}
   = {d \over df}\int\int{1 \over (y \log y)^{1 \over 2}}dy_m
   - \int C dx_m $$
   $$ = e^{-f} - e^{f} $$
   $$ {d \over df}F = {d \over df}\int \int{1 \over (x \log x)^2}dx_m 
   + {d \over df}\int\int{1 \over (y \log y)^{1 \over 2}}dy_m $$
   $$ = e^{f} + e^{-f} $$

     素粒子方程式は、２種類の素粒子と１種類の素粒子の組み合わせで、
   Jones多項式をペアーで構成されている。Rico level理論をベースに
   して、強い力と弱い力、電磁気力が、誤差をD-brane間を重力が
   gravitonとして伝わっているのが、洩れて、反重力をこの
   Weak electric theoryに加えると、全部の力が合わさって、
   統一場理論が完成される。
   これが、ヒッグス場とオイラーの定数を多様体を次元として、
   加群すると、ゼータ関数になる。これが、微分幾何の量子化とガンマ関数
   でもある。

   $$ ||ds^2|| = e^{-2 \pi T|\psi|}[\eta + \bar{h}(x)]dx^{\mu}dx^{\nu}
   + T^2 d^2 \psi $$
   
      素数分布論が、ゼータ関数を開集合として同型でもあり、ノルム空間、経路積分
   として、位相差が原子間力にもなっている。このエネルギー分布が、
   素数と素粒子方程式、宇宙と異次元のフェルミオンとして、一致している。
   このエネルギー分布が、複素多様体である。
   $$\square = 2 (\cos (i x \log x) + i \sin(i x \log x) $$ 


      宇宙と異次元が、双対性として、超対称性の素粒子を形成している。　


 $$ {\partial \over \partial f}F(x) = \int \int \text{cohom}D_k(x)[I_m]  $$
 Cohomology element is constructed with isotopy of D-brane,
 and this isotopy of symmetry structure is sheaf of manifold,
 more over this brane of element is double integrate of projection.
 $$ {\partial \over \partial f}F = {{}^t {\scalebox{1}[2]{{\variint}}}}
 \text{cohom}D_k(x)^{\ll p} $$
 And global partial defferential equation is integrate of cut in cohomolgy,
 and this step of defferential D-brane of sheap value is varint of equals,
 inverse of horizen of cute of section value is reverse of integrate of 
 operators. Three dimension of varint of manifolds in operator of sheaps,
 and this step of times of value is equal with D-brane of equations.
 Global differential equation and integrate equation is operator of 
 global scales of horizen cut and add of cut equation are routs.

 $$ \oint r dx_m = \mathcal{O}(x,y) $$

 $$ {{}^t {\scalebox{1}[2]{{\variiint}}}}_{{D(\chi,x)}} \text{Hom}[D^2 \psi]
 ^{\ll p} \cong  \text{vol}\left({V \over S}\right) $$

 $$ {\partial \over \partial f}F(x) = {{}^t \scalebox{1}[2]{\variint}} 
 \text{cohom}D_k(x)^{\ll p} $$

 $$ {d \over df}F = (F)^{f^{'}}, \int F dx_m = (F)^{f}  $$
 $$ {d \over df}\int \int{1 \over (x \log x)^2}dx_m =
 \left(\int \int{1 \over (x \log x)^2}dx_m \right)^{f^{'}}  $$
 $$ f^{'} = (2x (\log x + \log (x + 1))) $$
  $$  \int \int F dx_m = \left({1 \over (x \log x)^2}\right)^{
	  (\int 2x (\log x + \log (x + 1))dx)} $$ 
   $$= e^{-f} $$
  $$ {d \over df}F = \left({1 \over (x \log x)^2}\right)^{
	  (2x (\log x + \log (x + 1))) } $$  
  $$= e^{f} $$
  $$ \log (x \log x) \ge 2(y \log y)^{1 \over 2} $$

  $$ \log (x \log x) \ge 2(\sqrt{y \log y}) $$

  Global differential manifold is calucrate with x of x times in non
  entropy of defferential values, and out of range in differential formula
  be able to oneselves of global differential compute system.
  Global integrate manifold is also calucrate with step of resulted with
  loop of integrate systems. Global differential manifold is newton and 
  laiptitz formula of out of range in deprivate of compute systems.
  This system of calcurate formula is resulted for reverse of global equation
  each differential and intergrate in non entropy compute resulted values.

  $$ \pi(\chi,x) = [i\pi(\chi,x),f(x)] $$
  $$ \int{1 \over (x \log x)}dx = i \int{1 \over (x \log x)}dx^{N}  
  + {}^{N} i \int{1 \over (x \log x)}dx $$
  $$  \int{1 \over (x \log x)}dx = i \int{x \log x}dx 
  - \int{1 \over (x \log x)}dx $$
  $$ \int \int{1 \over (x \log x)^2}dx_m = i {1 \over 2}x^2 $$
  $$ \int \int{1 \over (x \log x)^2}dx_m = i \int \int_M dx_m $$
  $$ \le {1 \over 2}i + x^2 $$
  $$ E = - {1 \over 2}mv^2 + mc^2 $$
  $$ \lim_{x \to \infty} \int \int{1 \over (x \log x)^2}dx_m \ge {1 \over 2}i $$
  $$ {d \over df}\int \int{1 \over (x \log x)^2}dx_m = {1 \over 2}i $$
  $$ {d \over df}\int \int{1 \over (x \log x)^2}dx_m = 
  \left(\int \int {1 \over (x \log x)^2}dx_m \right)^{(f)^{'}}$$
  $$ (f)^{(f)^{'}} $$
  $$ = (f^f)^{'} $$
  $$ = e^{x \log x} $$

  $$ \Gamma = \int {e^{-x}x^{1-t}}dx, 
  {d \over df}F = e^f, \int F dx_m = e^{-f} $$
  $$ \int x^{1-t}dx = {d \over df}F, \int e^{-x}dx = \int F dx_m $$
  $$ e^{i \theta} = \cos \theta + i \sin \theta $$
  $$ x \bot y \to x \cdot y = \vec{0}, 
  {d \over df}\int \int{1 \over (x \log x)^2}dx_m = {1 \over 2}i $$

  $$ {d \over df}\int \int{1 \over (x \log x)^2}dx_m = 
  \left(\int \int {1 \over (x \log x)^2}dx_m \right)^{(f)^{'}}$$
  $$ {1 \over 2} + {1 \over 2}i = \vec{0}, {1 \over 2}\cdot{1 \over 2}i $$
  $$ A + B  =  \vec{0}, A \cdot B = {1 \over 4}i $$
  $$  \tan 90 \neq 0, ||ds^2|| = A + B $$
  Represented real pole is not only one of pole but also real pole of 
  $ {1 \over 2} $, 
  $\sin 0 = 0$ $$ e^{i \theta} = \cos \theta + i \sin \theta $$
  $\theta $
  this result equation is non-certain system concerned with particle, electric 
  particle and atoms in open set group with possiblity of surface values,
  and have abled to proof in this group of system. Universe and the other
  dimension of concerned of relationship values.
  $$ {d \over df}F = {d \over df} \int \int{1 \over (x \log x)^2}dx_m 
  + {d \over df} \int \int {1 \over (y \log y)^{1 \over 2}}dy_m
  = \Gamma(x,y) $$
  $$ {d \over d \gamma}\Gamma = (e^{-x}x^{1-t})^{\gamma^{'}} \to
  {d \over d \gamma}\Gamma = \Gamma^{(\gamma)^{'}}$$



    $$ \Gamma(s) = \int e^{-x} x^{s - 1}dx, \Gamma^{'}(s)
    = \int e^{-x} x^{s - 1}\log x dx $$
    $$ x^{s - 1} = e^{(s - 1)\log x} $$
    $$ {\partial \over \partial s}x^{s - 1} =
       {\partial \over \partial s}e^{-x}\cdot (e^{{s - 1}\log x}) 
       = e^{-x}\cdot \log x \cdot e^{(s - 1) \log x} $$
    $$ = e^{-x}\cdot \log x \cdot x^{s - 1} $$
    $$ {d \over d \gamma}\Gamma = \Gamma^{(\gamma)^{'}} $$
    $$ {d \over d \gamma}\Gamma(s) = \left( \int^{\infty}_{0}e^{-x}
    x^{s - 1}dx \right)^{\left(\int^{\infty}_{0}e^{-x}x^{s - 1}\log x dx
    \right)^{'}} $$
    $$ = \Gamma^{(\Gamma \int \log x dx)^{'}} $$
    $$ = e^{-x \log x} $$

    $$ \Gamma(s) = \int^{\infty}_{0} e^{-x}x^{s - 1}dx $$
    $$  \Gamma^{'}(s) = \int^{\infty}_{0} e^{-x}x^{s - 1} \log x dx $$
    $$ {d \over df}F = \int x^{s - 1}dx $$
    $$ \int F dx_m = \int e^{-x}dx $$
    $$ {d \over df}F = F^{(f)^{'}}, \int F dx_m = F^{(f)} $$
    Global differential manifold and global integrate manifold are
    reached with begin estimate of gamma and beta function leaded solved,
    Global differential manifold is emerged with gamma function of
    themselves of first value of global differential manifold of gamma
    function, this equation of quota in newton and dalarnverles of equation
    is also of same results equations. Zeta function also resulted with
    quantum group reach with quanto algebra equation.
        Included of diverse of gravity of influent in quantum physics,
    also quantum physics doesn't exist Gauss surface of theorem,
   
    $$ H \Psi = \bigoplus (i \hbar ^{\nabla})^{\oplus L} $$
    $$ = \bigoplus {{ H \Psi} \over \nabla L} $$
    $$ = e^{x \log x} = x^{(x)^{'}} $$
    however,
    $$ {d \over df}F = m(x) $$
    and gravity equation in being abled to existed with quantum physics,
    and this physics is also represented with quantum level of differential
    geometry.
    $$ {d \over dt}\psi(t) = \hbar $$
    $$ = {1 \over 2}i e^{i \Hat{H}} $$
    $$ (i \hbar)^{'} = (- e^{i \Hat{H}})^{'} $$
    $$ = - i e^{i \Hat{H}} $$
    $$ \psi(x) = e^{-i \Hat{H}t}, \bigoplus (i \hbar ^{\nabla})^{\oplus L} 
    = {{{1 \over 2} e^{i \Hat{H}}}^{(-i e^{i \Hat{H}})}} $$
    $$ = ({1 \over 2}f)^{-i f}$$
    $$ = ({1 \over 2})^{-i f} \cdot e^{-x \log x} \cdot (f)^{i} $$
    $$ = \int e^{- x} x^{t - 1}dx, {d \over d \gamma}\Gamma = e^{-x \log x}  $$
     
    $$ f = x, i = t, {1 \over 2} = a, \bigoplus a^{-t x}x^t [I_m]
     \cong \int e^{-x} x^{t - 1}dx $$
     Quantum level of differential geomerty is also resulted with
    Gamma function of global differential manifold of values, Heisenberg
    equation is alos resulted with global integrate manifold and this
    quantum level of differential geometry is manifold of deprive of formula.
    $$ |\psi(t)\rangle_s = e^{-i\Hat{H} t}|\Psi\rangle_{H},
    \Hat{A}_s = \Hat{A}_H(0) $$
    $$ |\Psi(t)\rangle_s \to {d \over dt} $$
    $$ i {d \over dt}|\psi (t) \rangle_s = \Hat{H}|\psi(t)\rangle_s $$
    $$ \langle\Hat{A}(t)\rangle = \langle\Psi(t)|\Hat{A}(0)|\Psi(t)\rangle  $$
    $$ {d \over dt}\Hat{A} = {1 \over i}[\Hat{A},H] $$
    $$ \Hat{A}(t) = e^{i\Hat{H}t}\Hat{A}(0)e^{-i\Hat{H}t} $$
    $$ \lim_{\theta \to 0} \begin{pmatrix} \sin \theta \\
    \cos \theta \end{pmatrix} \begin{pmatrix} \theta & 1 \\
    1 & \theta \end{pmatrix} \begin{pmatrix} \cos \theta \\
    \sin \theta \end{pmatrix} = \begin{pmatrix} 1 & 0 \\
    0 & - 1 \end{pmatrix} $$
    $$ f^{-1}(x) x f(x) = I^{'}_m, I^{'}_m = [1,0] \times [0,1] $$

    $$ {x + y} \ge \sqrt{x y} $$
    $$ {x^{{1 \over 2} + iy} \over e^{x \log x}} = 1 $$
    $$ \mathcal{O}(x) = \nabla_i \nabla_j \int e^{{2 \over m}\sin \theta 
    \cos \theta } \times {{N \mathrm{mod}(e^{x \log x})} \over {O(x)(x +
    \Delta |f|^2)^{1 \over 2}}} $$
    $$ x \Gamma(x) = 2 \int|\sin 2 \theta|^2 d \theta, \mathcal{O}(x)
    = m(x)[D^2 \psi] $$
    $$ i^2 = (0,1) \cdot (0,1), |a||b|\cos \theta = - 1 $$
    $$ E = \mathrm{div}(E, E_1) $$
    $$ \left({\{f,g\}} \over [f,g] \right) = i^2, E = mc^2, I^{'} = i^2 $$
     Gamma function of global differential equation is first of differential
    function deprivate with ordinary differential equations, more also
    universe and other dimensiion of Higgs field emerge with zeta function,
    more spectrum focus is three dimension of manifold of singularity or pair
    theorem, and universe and other dimension of each value equals.
    This reverse of circle function of manifolds is also Higgs fields of 
    dense of entropy and result equation. And differential structure of quantum
    level in gravity is poled with zeta function and quantum group in high 
    result values. Zeta funtion and quantum group is Global differential
    manifold of result in values.
    $$ {d \over d \gamma} \Gamma = m(x) $$
    $$ = e^{- x \log x} $$
    $$ \sin i x = {{e^{-x} + e^{x}} \over 2i} $$
    $$ {d \over df}F = m(x) = e^{f} + e^{-f} $$
    $$ = 2 i \sin (i {x \log x}) $$

    Hyper circle function is constructed with sheap of manifold in 
    Beta function of reverse system emerged from Gamma fucntion 
    and reverse of circle function of entropy value,
    and this reverse of hyper circle function of beta system
    is universe and the other dimension of zeta function,
    more over spectrum focus is this beta function of reverse in
    circle function is quantum level of differential structure oneselves,
    and this quantum level is also gamma function themselves.
    This gamma fucntion is more also D-brane and anti-D-brane built 
    with supplicant values. Higgs fileds is more also pair of universe and
    other dimension each other value elements. 


    My son in acasicrecord demonstrate with being have with quantum level
    of differential structure equation.
    I have with same equation already from being on 
    my father and my doctors of professors give me hints with 
    this qlds equation explanade with global differential manifold.

    $$ H\Psi = \bigoplus {(i \hbar^{\nabla})}^{\oplus L}  (1)$$
    $$ = i\hbar\psi $$
    $$ {d \over df}F = m(x) (2) $$
    $$ = {d \over df}\int \int{1 \over (x \log x)^2}dx_m
    + {d \over df}\int \int{1 \over (y \log y)^{1 \over 2}}dy_m  (3)$$
    $$  \bigoplus {(i \hbar^{\nabla})}^{\oplus L} = \int e^{-x}x^{1 - t}dx_m 
    (4)$$
    $$ = {d \over d \gamma}\Gamma = \int \Gamma(\gamma)^{'}dx_m (5)$$
    $$ = e^f + e^{-f} (6)$$
    $$ {d \over d \gamma}\Gamma^{-1} = e^{f} - e^{-f} (7)$$
    (1) is my son in acasicrecord with his idea, (2),(3) are myself equation.
    (4),(5) are my son and me with tagge of ideas equation.
    (6),(7) are Takashima Aya and me, and my son in acasicrecode with tolio
    equations.
    These references from Grisha professor and Takeuchi Kaoru.
    I read with this references being hint from these professors.
 
      [reference Lisa Randall, Toshihide Masuoka, Makoto Kobayashi]



   [Reference: 深谷賢治先生と加藤文元先生と竹内薫先生、リサ・ランドール博士、
   S.マックレーン先生]

\title{Inspectable theorem in rotate of each dimension \\
symmetry broken in world tree emerged with \\
energy of entropy non commutative from \\
formation of live in tree waters.}


        Adam and eve are eaten with apple of live in tree for acasic record.
	Therefore, A apple were eaten from Adam and eve firstly.
	Each distance of experiences in same happend events.


        Entropy formula demonstrate with zeta function destrout from all of
    world tree in stem of groundlines. This stem of function instrumentate
    for many retree of being endevoured into water of sea in ocean.
    These norm of being rested with forever of deepist learning,
    and this read with book of acasicrecord formate with live of water,
    or this destrate with water into black hole and center of line
    in white hole, till of center of broken in atom levels,
    stephaniy of atom received with quarks of particle level
    for CP violation in AdS5 manifold on universe and other dimension.
    
\title{特殊相対性理論の虚数回転による多様体積分と、\\
それによる一般相対性理論の再構築理論}
\author{}
   
  \newcommand{\oiint}{%
	   \begin{picture}(15,15)
		   \put(0,0){
			   $ \scalebox{1}[2]{$\iint$} $}
		   \put(2,2){$\bigcirc$}
	   \end{picture} %
	   }


 $$ {\partial \over \partial f}F(x) = \int \int \text{cohom}D_k(x)[I_m]  $$
 Cohomology element is constructed with isotopy of D-brane,
 and this isotopy of symmetry structure is sheaf of manifold,
 more over this brane of element is double integrate of projection.
 $$ {\partial \over \partial f}F = {{}^t \scalebox{1}[2]{\variint}}
 \text{cohom}D_k(x)^{\ll p} $$

  $$  {\partial \over \partial f}F = {{}^t \scalebox{1}[2]{\variint}}
 \text{cohom}D_k(x)^{\ll p} = 
\scalebox{1}[2]{\alearletter}$$ 
   $$ \nabla_{i} \nabla_{j} \int f(x) d \eta = {\partial^2 \over {\partial x
   \partial y}}\int \scalebox{1}[2]{\alearletter} d \eta $$
   一般相対性理論の加群分解が偏微分方程式と同じく、
   特殊相対性理論の多様体積分の虚数回転体がベータ関数となる。
   ほとんどの回転体の体積が、係数と冪乗での回転体として、ベータ関数
   と言える。
   $$ = R^{{\mu\nu}^{'}} + {1 \over 2}\Lambda g_{ij}^{'} = \int 
   \left(i {{v \over {
	   \sqrt{1 - 
   {{{({v \over t})}}^{2}}}}}} + {v \over {\sqrt{1 - 
   {{{({v \over t})}}^{2}}}}}\right)\mathrm{dvol} $$
   $$ = {d \over df}\int \int{1 \over {(x \log x)^2}}dx_m +
   {d \over df}\int \int {1 \over {(y \log y)^{1 \over 2}}}dy_m $$ 
    この大域的積分多様体が大域的微分多様体の反重力と重力方程式
    で表せられて、
   $$ = \int C dx_m = \int \kappa T^{\mu\nu} \mathrm dx_m = T^{{\mu\nu}{^{T^
   {\mu\nu}}}^{'}}$$
   オイラーの定数の大域的積分多様体が、一般相対性理論の大域的積分多様体
   であり、エントロピー不変式で表せられる。

  $$ \scalebox{1}[2]{\alearletter} = 
  {{{{8 \pi G} \over {c^4}}T^{\mu\nu}}/ \log x} $$
  $$ {}^{t}\scalebox{1}[2]{\variiint}\mathrm{cohom D_{\chi}[I_m]}$$
  $$ = \oint {(px^n + qx + r)^{\nabla l}} $$
  $$ {d \over d l}L(x,y) = 2 \int ||\sin 2x||^{2}d \tau $$
  $$ {d \over d \gamma}\Gamma $$
 $$ ||ds^2|| = \lim_{x \to \infty} [\delta(x) \int \int \int \pi \left( \sum_{k=0}^{\infty}
		{{}^n\sqrt{p},x \over n} \right)
^{1 \over 2}d \tau]^{\mu\nu} $$

 $$\oiint  \cong  ||ds^2|| = \lim_{x \to \infty} 
 [\delta(x) \int \int \int \pi \left( \sum_{k=0}^{\infty}
		{{}^n\sqrt{p},x \over n} 
		\right)^{1 \over 2}d \tau]^{\mu\nu} $$
 
 $$ e^{-2\pi T||\psi||}[\eta + \bar{h}]dx^{\mu}dx^{\nu} + T^2 d^2 \psi $$
 $$ {{-16 \pi G} \over c^4}T^{\mu\nu}/ \log x =  
 \scalebox{1}[2]{\alearletter} $$
 $${{-16 \pi G} \over c^4}T^{\mu\nu}/ e^{-2\pi T||\psi||} $$
 $$ = 4 \pi G \rho $$
 $$ {\partial \over {\partial x \partial y}} = \nabla_{i} \nabla_{j} $$ 
 $$ \square \iiint  = {}^{t}\scalebox{1}[2]{\variiint}\ $$
 $$ {\partial \over {\partial x}} \iiint = \nabla_{i}\nabla_{j} \int
 \nabla f(x) d \eta $$
 $$ {\partial \over {\partial f}}\iiint = \square \iiint $$

 $$ \reflectbox{$\Gamma$} | \Gamma, \reflectbox{B} | B $$
 $$ \reflectbox{E} | E, \reflectbox{C} | C $$
 $$ \reflectbox{F} | F, \reflectbox{$\beta$} | \beta $$
 $$ \reflectbox{D} | D, {}^{t}\scalebox{1}[2]{\variint} \cong 
 \bigoplus D^{\bigoplus L} $$
 $$ \square +  \scalebox{1}[2]{\alearletter} = \emptyset $$
 $$  \scalebox{1}[2]{\alearletter} | \emptyset = \square $$
 $$ {\begin{pmatrix} 1 & 0 \\
 0 & -i \end{pmatrix}}^{1 \over 2} \Biggl| {\begin{pmatrix} 1 & 0 \\
 0 & -i \end{pmatrix}}$$
 $$ x^{1 \over 2} \cong x, \emptyset^{1 \over 2} = \emptyset$$
 
    自分がいると、方位が決まる。そうしたら、空間ができる。北と南が出来る。
    これは、磁気単極子が相対空間での根本的なアイデアらしい。
    無になると関係性がなくなり、磁気単極子もある上に、全て無でもある。
    無になると十二支縁起になる。宇宙と異次元の外が無でもある。
    運動量が決まると位置が決まる。位置だけだと宇宙が不確定性原理
    であり、異次元が合わさると運動量も位置も決まる。
    無になると位置も運動量も同時にわかる。
    このアイデアが3次元多様体のエントロピー値である。
    
    $$ \int^{\ll n} ||ds^2|| = \int\int\int 8 \pi G \left({p \over c^3} + 
    {V \over S}\right)dxdydz $$
    $$ = {\partial \over \partial f_n}{||ds^2||} $$


 $$ {\partial \over \partial f}F(x) = \int \int \text{cohom}D_k(x)[I_m]  $$
 Cohomology element is constructed with isotopy of D-brane,
 and this isotopy of symmetry structure is sheaf of manifold,
 more over this brane of element is double integrate of projection.
 $$ {\partial \over \partial f}F = {{}^t \scalebox{1}[2]{\variint}}
 \text{cohom}D_k(x)^{\ll p} $$
 And global partial defferential equation is integrate of cut in cohomolgy,
 and this step of defferential D-brane of sheap value is varint of equals,
 inverse of horizen of cute of section value is reverse of integrate of 
 operators. Three dimension of varint of manifolds in operator of sheaps,
 and this step of times of value is equal with D-brane of equations.
 Global differential equation and integrate equation is operator of 
 global scales of horizen cut and add of cut equation are routs.

 $$ \oint r dx_m = \mathcal{O}(x,y) $$

 $$ {{}^t \scalebox{1}[2]{\variiint}}_{{D(\chi,x)}} 
 \text{Hom}[D^2 \psi]^{\ll p} \cong 
 \text{vol}\left({V \over S}\right) $$

 $$ {\partial \over \partial f}F(x) = {{}^t \scalebox{1}[2]{\variint}} 
 \text{cohom}D_k(x)^{\ll p} $$

 $$ {d \over df}F = (F)^{f^{'}}, \int F dx_m = (F)^{f}  $$
 $$ {d \over df}\int \int{1 \over (x \log x)^2}dx_m =
 \left(\int \int{1 \over (x \log x)^2}dx_m \right)^{f^{'}}  $$
 $$ f^{'} = (2x (\log x + \log (x + 1))) $$
  $$  \int \int F dx_m = \left({1 \over (x \log x)^2}\right)^{
	  (\int 2x (\log x + \log (x + 1))dx)} $$ 
   $$= e^{-f} $$
  $$ {d \over df}F = \left({1 \over (x \log x)^2}\right)^{
	  (2x (\log x + \log (x + 1))) } $$  
  $$= e^{f} $$
  $$ \log (x \log x) \ge 2(y \log y)^{1 \over 2} $$

  $$ \log (x \log x) \ge 2(\sqrt{y \log y}) $$

 
    $$ ||ds^2|| = 8 \pi G \left({p \over c^3} + {V \over S}\right) $$
    $$ \int ||ds^2||dx_m  = \int 8 \pi G \left({p \over c^3} 
    + {V \over S}\right)\mathrm{dvol}  $$
    $$ \int ||ds^2||dx_m  = \int 8 \pi G \left({p \over c^3} 
    + {V \over S}\right)dy_m $$
     $$ \int ||ds^2||dx_m  = \int {1 \over {8 \pi G \left({p \over c^3} 
    + {V \over S}\right)}}dx_m $$
    $$ {d \over df}F = m(x), \bigoplus \left({i \hbar}^{\nabla}\right)^{
	    \oplus L} $$
    $$ = \nabla_{i} \nabla_{j} \int \nabla f(x)d \eta $$

    $$ dx_m = {y \over {\log x}}, dy_m = {x \over \log x} $$
    $$ e^{-f}dV = dy_m = \mathrm{dvol} $$
    偏微分は加群分解と同じ計算式に行き着く。

   宇宙と異次元の誤差関数のエネルギー、
    $AdS_5$多様体がベータ関数となる値の列が、異次元への扉となっている。
    $$ \beta(p,q) = \text{誤差関数} + \text{Abel多様体} $$
    $$ = \text{$AdS_5$多様体} $$
    $$ = {d \over df}F + \int C dx_m = \int \Gamma(\gamma)^{'}dx_m $$
    ここで、アーベル多様体はEuler productである。
    ベータ関数の数列がわかると、ゼータ関数は無であるというのが、
    どういうことかが、物、物体に影ができて、ものが瞑想と同じであり、
    これから、風景がベータ関数の数値列に見えるらしい。
     この大域的微分多様体のガンマ関数が、複素力学系のマンデルブロ集合の
    プリズムと同じ構造の見方らしい。

    $$ ||ds^2|| = e^{-2\pi G ||\psi||}[\eta + \bar{h}(x)]dx^{\mu}dx^{\nu}
    + T^2 d^2 \psi $$
    $$ \beta(p,q) = \text{誤差関数} + \text{Abel多様体} $$
    $$ \int \mathrm{dvol} = \square \psi $$
    $$ \int \nabla \psi^2 d \nabla \psi = \square \psi $$
    $$ \text{expanding of universe} = \text{exist of value} $$
    $$ = \log (x \log x) = \square \psi $$
    $$ \text{freeze out of universe} = \text{reality of value} $$
    $$ = (y \log y)^{1 \over 2} = \nabla \psi $$

    All of value is constance of entropy, universe is freeze out
    constant, and other dimension is expanding into fifth dimension of
    inner.


      $$ x^n + y^n = z^n, \beta(p,q) = x^n + y^n - \delta(x) = z^n - \delta
      (x) $$
      $$ {d \over dt}g_{ij} = - 2 R_{ij}, {\cos x \over (\cos x)^{'}}
      \cdot (\sin x)^{'} = z_n, z^n = - 2 e^{x \log x} $$

      $$ z = e^{-f} + e^{f} - y $$
      $$ \beta(p,q) = e^{-f} + e^{f} $$
      相対性は、暗号解読と同じ仕組みの数式を表している。
      ここで言うと、yが暗号値である。テェックディジットと同じ仕組みを有
      している。

      $$ \square x = \int {f(x) \over {\nabla(R^{+}\cap E^{+})}}d \square x$$
      $$= \int {{\Delta f(x) \circ E^{+}} \over {\nabla(R^{+}\cap E^{+})}}
      \square x $$
      $$ \square x = \int {d(R \nabla E^{+}) \over {\nabla(R^{+} \cap E^{+})}}
      d \square x $$
      $$ = \int {{\nabla_{i} \nabla_{j} (R + E^{+}} \over {\nabla(R^{+}
      \cap  E^{+})}}d \square x $$
      $$ x^n + y^n = z^n $$
      $$ \mathrm{exp}(\nabla(R^{+}\cap E^{+}),\Delta(C \supset R)) $$
      $$ = \pi (R_1 \subset \nabla E^{+}) = \mathrm{rot}(E_1, 
      \mathrm{div}E_2) $$
      $$ xf(x) = F(x), s\Gamma(s) = \Gamma(s+1) $$
      $$ Q \nabla C^{+} = {d \over df}F(x) \nabla \int \delta(s)f(x)dx $$
      $$ E^{+}\nabla f = {{e^{x \log x} \nabla n! f(x)} \over {E(x)}} $$
      $$ {d \over df}F ~ F^{{f}^{'}} = e^{x \log x} $$
      $$ (C^{\nabla})^{\oplus Q} = e^{x \log x} $$
      $$ {d \over dt}g_{ij} = - 2 R_{ij}, {\cos x \over (\cos x)^{'}}
      \cdot (\sin x)^{'} = z_n, z^n = - 2 e^{x \log x} $$

      $$ R \nabla E^{+} = f(x) \nabla e^{x \log x}, {d \over df}F = F^{
	      {f}^{'}} = e^{x \log x}  $$
      南半球と単体（実数）の共通集合の偏微分した変数をどのようなF(x)かを
      $$ \int \delta(x)f(x)dx $$
      と同じく、単体積分した積分、共通集合の偏微分をどのくらいの微分変数
      を
      $$ \int \nabla \psi^2 d \nabla \psi = \square \psi $$
      と同じ、
      $$ \int dx = x + \text{C(Cは積分定数) }$$
      と原理は同じである。


      Beta function is,
      $$ \beta(p,q) = \int x^{1 - t} (1 - x)^{t}dx = \int t^{x}(1 - t)^{x - 1}
      dt $$
      $$ 0 \le y \le 1, \int^{1}_{0} x^{10}(1 - x)^{20}dx = B(11,21) $$
      $$ = {{\Gamma(11)\Gamma(21)} \over \Gamma(32)} = {{10! 20!} \over {
	      31!}} = {1 \over {931395465}} $$
       $$ {1 \over {931395465}} = {1 \over 9} = {1 \over {1 - x}} $$
	 $$     = {1 \over {1 - z}} =
	      \sum_{k=0}^{\infty}z^k = {1 \over {1 + z^2}} = \sum_{
		      k=0}^{\infty} (-1)^k z^{2k} $$
      $$ f(x) = \sum_{k=0}^{\infty} a_k z^k $$
      $$ {{d^n y} \over dx^n} = n! y^{n+1} $$
      $$ f^{(0)}(0) = n!f(0)^{n+1} = n! $$
      $$ f(x) \cong \sum_{k=0}^{\infty}x^n = {1 \over {1 - x}} $$
      In example script is,
      $$ {{dy} \over dx} = y^2, {1 \over y^2}\cdot {dy \over dx} = 1 $$
      $$ \int {1 \over y^2}{dy \over dx}dx = \int{dy \over y^2} = - {1 \over y
      } $$
      $$ - {1 \over y} = x - C, y = {1 \over {C - x}} $$
      $$ \exists x = 0, y = 1 $$
      $$\text{in first value condition compute with} $$
      $$ \text{result, C consumer sartified,} $$
      $$ y = {1 \over {1 - x}} $$
      This value result is concluded with native function from Abel manifold.

      Abelian enterstane with native funtion meltrage from Daranvelcian
      equation on Beta function.
 
       Rotate with pole of dimension is converted from 
       other sequence of element,this atmosphere of vacume 
       from being emerged with quarks of being constructed
with Higgs field, and this created from energy of zero dimension are begun with
universe and other dimension started from darkmatter that represented with 
big-ban.
Therefore, this element of circumstance have with quarks of twelve lake with
atoms.

  $$ (dx, \partial x)\cdot (\epsilon x, \delta x) = \begin{pmatrix}
	  1 & 0 \\
  0 & - 1 \end{pmatrix}^{1 \over 2}  $$
  $$ {d \over df} F = m(x) $$

 And this phenomounen is super symmetry theorem built with quarks of element, 
also this chemistry of mechanism estimate from physics of operatsion.
These response of mechanism chain of geometry into space being emerged with
creature and univse of existing of combination.

  This converted with dimension of element also emerged from imaginary and 
reality of pole's space.

  And this pole of transport of dimension belong with vector of time has 
with one rout of sequecence. Quanum physics also belong with other vector 
of time has with imaginary rout of sequecense.

 This sequence of being estimated with non fluer of time, and this space 
of element have with gravity and antigravity of power.
Other vecor of time is antigravity rout of sequecense.

 Weak electric theory is estimate from time has with one rout of sequecense,
this topology of chain is resulted from time of rout ways.

      $$ \square ({{\sigma_1 + \sigma_2} \over 2 }) = [3 \pi (\chi,x) \circ f(x),
   \sigma(x)] $$
   $$ \square \psi = 8 \pi G T^{\mu\nu} $$
   $$ {d \over dt}g_{ij}(t) = - 2 R_{ij} $$
   $$ \nabla \psi^2 = 4 \pi G \rho $$
   $$ \square (\sigma_1 + \sigma_2) = [6\pi(\chi,x)\circ f(x), \sigma(x)] $$
   $$ = [i \pi(\chi,x),f(x)] $$
   $$ \square \psi = [12 \pi (\chi,x)\circ f(x), \sigma(x)] $$
   $$ {d \over df} F = \int e^{-f}[- \Delta v + R_{ij} v_{ij}+ 
   \nabla_i \nabla_j f  + v \nabla_i \nabla_j + 2<f,h> + (R + \nabla f)
   ({v \over 2} - h) ] $$
   $$ =[i \pi(\chi,x), f(x)] $$

   These equation is reminded time pass rout of one rout way of forms,
   and this rout of time ways which go for system from future and past.
   Therefore, this resulted system of time mechanism is one true flow
   that weak electric theorem oneselves.
   and moreover, one rout time way of forms is reverse with antigravity of
   time system. This also spectrum focus is true that Maxwell theorem and
   strong boson unite with antigravity, this unite is essense on the contrary
   from weak electric theorem, this theorem called for time rout forms is
   strong electric theorem.
   This two theorem is united with quantum physics that no time flow system.

   $$ f^{-1}(x)xf(x) = 1,H_m = E_m \times K_m $$


   The non-commutative theorem is constructed from world line surface that
   this complex manifold estimate with rolanz atractor, and this string
   theorem have with one world of universe mate six quarks and other world of
   dimension mate other element of six quarks.
   These particle quarks is built with super symmetry space of dimension.

   $$ i = (1,0)\cdot (1,0), e^{i\theta} = \cos \theta + i \sin \theta $$
   $$ e^{-i \theta} = \cos \theta - i\sin \theta $$
   $$ \sin \theta = {e^{i \theta} - e^{-i \theta} \over 2i} $$
   $$ \sin i \theta = {e^{-\theta} - e^{\theta} \over 2} $$
   $$ \pi (\chi,x) = \cos \theta + i \sin \theta $$

   In this equations, two dimension redestructed into three dimension,
   this destroy of reconstructed way is append with fifth dimension.
   This deconstructed way of redestructe is arround of universe attached
   with three dimension, this over cover call into fifth dimension.

   $$ R(-\alpha) = \begin{pmatrix} \cos \alpha & \sin \alpha \\
   - \sin \alpha & \cos \alpha \end{pmatrix} $$
   $$ R(\alpha) = \begin{pmatrix} \cos \alpha & - \sin \alpha \\
   \sin \alpha & \cos \alpha \end{pmatrix} $$
   $$ R(\alpha)M R(-\alpha) = 
   \begin{pmatrix} \cos \alpha & - \sin \alpha \\
   \sin \alpha & \cos \alpha \end{pmatrix} \begin{pmatrix} 1 & 0 \\
   0 & - 1 \end{pmatrix} \begin{pmatrix} \cos \alpha & \sin \alpha \\
   - \sin \alpha & \cos \alpha \end{pmatrix} $$
   $$ = \begin{pmatrix} \cos^2 \alpha - \sin^2 \alpha & 2 \sin \alpha \cos 
	   \alpha \\
	   2 \sin \alpha \cos \alpha & - \cos^2 \alpha + \sin^2 \alpha
   \end{pmatrix} $$
   $$ = \begin{pmatrix} \cos 2 \alpha & \sin 2 \alpha \\
   \sin 2 \alpha & - \cos 2 \alpha \end{pmatrix} $$

   $$ \begin{pmatrix} \cos \alpha & - 1 \\
	   1 & - \sin \alpha \end{pmatrix} \begin{pmatrix} \cos \alpha \\ 
		   \sin \alpha
   \end{pmatrix} = \begin{pmatrix} 1 & 0 \\
   0 & - 1 \end{pmatrix} $$
   $$ \begin{pmatrix} \cos \alpha & \sin \alpha \\
   \sin \alpha & - \cos \alpha \end{pmatrix} \begin{pmatrix}
   x \\ y \end{pmatrix} = \begin{pmatrix} 1 & 0 \\
   0 & - 1 \end{pmatrix} $$
   $$ \lim_{\theta \to 0} \begin{pmatrix} \sin \theta \\ \cos \theta 
   \end{pmatrix} \begin{pmatrix} \theta & 1 \\
   1 & \theta \end{pmatrix} \begin{pmatrix} \cos \theta \\ \sin \theta
	   \end{pmatrix} = \begin{pmatrix} 1 & 0 \\ 0 & - 1 \end{pmatrix} $$
		   $$ f^{-1}xf(x) = 1 $$
		   $$ (\log x)^{'} = {1 \over x}, x^n + y^n = z^n  $$
		   $$ x^n = - y^n + c, nx^{n-1} = -n y^{n-1}y^{'} $$
		   $$ y^{'} = {{n x^{n-1}} \over {n y^{n-1}}} $$
		   $$ = {x^{n-1} \over y^{n-1}}= - {y \over x}\cdot
		   ({x \over y})^n $$
		   $$ - {\cos x \over (\cos x)^{'}}(\sin x)^{'} = z_n $$
		   $$ z^n = - 2 e^{x \log x} $$
		   $$ \lim_{x \to \infty} f(x) = a, \lim_{y \to \infty}
		   f(y) = b, \lim_{x,y \to \infty}\{f(x) + f(y)\} = a + b  $$
		   $$ \lim_{z \to 0}f(z) = c, \delta \int z^n = {d \over dV}
		   x^3 $$
		   $$ \lim_{x \to \infty} = c - \lim_{y \to \infty}f(y)
		   \lim_{x \to \infty} f(x) + \lim_{y \to \infty}f(y) = 
		   \lim_{z \to \infty}f(z) $$
		   $$ z^n = \cos n \theta + i \sin n \theta $$
		   $$ = - 2 e^{x \log x} $$

  These equation is gravity and antigravity equation. 

	   $$ {d \over d\sigma}\begin{bmatrix}{(\sigma_1 + \sigma_2) 
	   \over 2}\end{bmatrix}  $$
	   $$ = \sigma(\upharpoonleft \downharpoonleft) + \sigma(
	   \upuparrows)+ \sigma(\downdownarrows)+ 
	   \sigma(\leftleftarrows)
	   + \sigma(\rightrightarrows) $$
	   $$ {d \over df}{\int \int {1 \over (x \log x)^2}}dx_m
	   = \sigma(\upharpoonright) $$
	   $$ {d \over df}{\int \int {1 \over (y \log y)^{1 \over2}}}dy_m
	   = \sigma(\downharpoonright) $$
	   $$ \sigma(\leftarrowtail) + \sigma(\rightarrowtail)
	   = \int e^{-f}[-\Delta v + R_{ij}v_{ij} + \nabla_{i}\nabla_{j}v
	   + v\nabla_{i}\nabla_{j} + 2 <f,h> + (R + \nabla f)(v - 
	   {h \over 2})]
	   $$
	   $$ \sigma(\downharpoonleft) = \sigma(\upharpoonright 
	   \downharpoonright + \upuparrows 
	   + \rightrightarrows)$$
	   $$ \sigma(\upharpoonleft) =  \sigma(\upharpoonright 
	   \downharpoonright + \downdownarrows 
	   + \leftleftarrows) $$
         
	   $$ \text{weak electric theorem} = \sigma(\upharpoonleft) $$
	   $$ \text{strong electric theorem} = \sigma(\downharpoonleft) $$

	   These equation is represented with topology of string model,
	   and weak electric theorem is constructed with Maxwell theorem 
	   and weak boson, gravity that estimate with this three power united.
	   Moreover, strong boson and Maxwell theorem, antigravity that
	   also estimate with this three power united.
	   This two united power is integrated from gravity and antigravity.
	   Then this united power is zeta function.

  Euler equation represet form,
  $$ C = \int {1 \over x^s}dx - \log x $$
  $$ = \int \left(\int {1 \over x^s}dx - \log x \right)\mathrm{dvol} $$
  $$ = \int \int{1 \over x^s}\mathrm{dvol} - \int  \log x \mathrm{dvol} $$
  This equation componet of $ \text{Higgs equation + Euler equation =
  Zeta function} $. This represent expression is Rich flow equation.
  This equation assent from Differential geometry in quantum level,
  Gamma and Beta function, Gauss surface, Higgs field, Einsetin tensor,
  Euler constance and so on.
    This equation is add and average of possibility, and universe, atom
  more over a certain theory, this universe oneselves is environment
  seek, other distance is possibility equation, however, other dimension
  cover demand with not only norm of non liner in three manifold
  but also non possibility in complete of environment equation.
    $$ = {d \over df} \int \int{1 \over (y \log y)^{1 \over 2}}dy_m 
  - \int e^{-x}x^{1 - t} \log x dx_m $$
  $$ \int x^{1 - t}e^{-x}dV = \int x^{1 - t}dm, \int x^{1 - t}e^{-x}dV
  = \int x^{1 - t}\mathrm{dvol}, f = \gamma = \Gamma^{'} = \int 
  e^{-x}x^{1 - t} \log x dx $$
  $$ = {d \over d \gamma}\Gamma^{-1} - (\gamma)^{\gamma^{'}} $$
  $$ = e^{-f} - e^{f} $$
  $$ = {d \over df}\int \int{1 \over (y \log y)^{1 \over 2}}dy_m
  - {d \over df}\int\int {1 \over (x \log x)^2}dx_m $$
  $$ = 2 \cos (i x \log x) $$
  $$  {d \over df}\int\int{1 \over (y \log y)^{1 \over 2}}dy_m
  + {d \over df}\int\int{1 \over (x \log x)^2}dx_m $$
  $$ = {d \over df}F = 2 i \sin (i x \log x) $$
  $$ {d \over df}F + \int C dx_m  
  = 2 (\cos( i x \log x) + i \sin(i x \log x)) $$
  $$ = 2 e^{- \theta} = 2 e^{- i x \log x} $$
  $$ = {d \over dt}g_{ij}(t) = - 2 R_{ij} $$
  $$ \log (x \log x) \ge 2 \sqrt{y \log y} $$
  $$ \log(x \log x) = \log x + \log \log x $$
  This universe and other dimension on blackhole and whitehole
  compont with universe and atom in movement and point of 
  environment value by a certain theory of complete environment 
  equation. And this three manifold of possibility equation 
  are able to research in both value formula.
  This envirnment of researchable reason is that surround of universe
  is setting to value by already in consented to being researched,
  then both value researchabled.
  And this three manifold equation is insert with Kaluza-Klein theory
  and Heisenbelg equation.
    $$ \nabla \psi^2 = 8 \pi G\left({p \over c^3} + {V \over S}\right) $$
  $$ \nabla \psi^2 = 8 \pi G \hbar + 8 \pi {V \over S} $$
  This energe of entropy value consist with whitehole of atom
  and universe in blackhole value, and this universe in balckhole 
  represent with Higgs filed, atom in whitehole consist with plack scale.
    $$ \square_v = 2 \sqrt{2 \pi G \hbar} + 2 \sqrt{2 \pi {V \over S}} $$
  This entropy value also consist with atom in balckole.
    $$ \sqrt{2 \pi T} = 2 \sqrt{2 \pi{V \over S}} $$

  These equation more also consist with universe being non gravity formula.
    $$ 2 \cos (i x \log x) + 2 i \sin(i x \log x) = 2 e^{-f} $$
  $$ 2 \sqrt{2 \pi {V \over S}} = {d \over df} 
  \int \int{1 \over (x \log x)}dx_m$$
  And also these equation consist with universe and other dimension in
  warped passeage in rout of integral.

   $$ e^{-2\pi ||psi|}[\eta_{\mu\nu} + \hbar{x}]dx^{\mu}dx^{\nu} 
  +  T^2 d^2 \psi $$
  $$ = e^{-f} + e^f = (u + d) + c, (w + s) + b = -e^{-f} + e^f $$
  この式は、時空にどれだけの原子が凝縮しているかを表していて、
  その上に、この逆数は、密度エネルギーをも示している。
  空間の体積にに原子を商代数にすると、濃度にもなる。
   逆数は、$\text{a=原子のエネルギー}$が全空間を1とすると
  $\text{$1 \over {y \over 
  x}$}$すると、$x \over y $は密度になる。　
  $\text{{x=原子},y=全空間、個数は}$$\text{$y \over x$}$ 
  $$ \square = \text{原子} \to {{[\eta_{\mu\nu} + \hbar(x)]dx^{\mu}dx^{\nu}
  } \over {e^{2 \pi T ||\psi||} }}$$
  $$ ||ds^2|| = e^{2 \pi T||\psi||}\left([\eta_{\mu\nu} + \hbar(x)]d^{\mu}dx^
  {\nu}\right)^{-1} + \left(T^2 d^2 \psi \right)^{-1} $$
  $$||ds^2|| = e^{-2 \pi T||\psi||}[\eta_{\mu\nu} + \hbar(x)]d^{\mu}dx^
  {\nu}] + T^2 d^2 \psi  $$
  $$ (u + d) + c = e^{-f} + e^f, (w + s) + b = e^{-f} + e^{f} $$
  $$ R_{ij}^{'} = - R_{ij} $$

  この$ AdS_5 $多様体の $e^{-2 \pi T ||\psi||} $は原子間距離を表していて、
  $ [\eta + \bar{h}(x)]d^{\mu}dx^{\nu} $で宇宙のD-braneの構造を表している。
  この数式が$T^2 d^2 \psi $とアーベル多様体が包み込んでいる。
  原子が回転するのと、ニュートンリングが回転する宇宙は同じ速さで回転する。　

  $$||ds^2|| = \square + \rho , ||ds^2|| = \nabla \Psi^2 + \psi $$
  $$ \eta_{\mu\nu} + \bar{h}(x) $$
  $$ \text{ploximetry} = \Delta x \Delta p - \Delta p \Delta x + 
  \delta(p,x)  $$
  $$ ||ds^2|| = e^{-2 \pi T ||\psi||} + [\eta_{\mu\nu}  + \bar{h}(x)]
  dx^{\mu\nu} + T^2 d^2 \psi $$
  $$ {d \over df}F = e^f + e^{-f}, {d \over df}\int C dx_m = e^f - e^{-f} $$
  $$ R_{ij}^{'} = - R_{ij} $$
  宇宙と異次元では０だが宇宙だけだと誤差が生じる。これが不確定性原理であり、
  異次元が誤差になり、宇宙と加群すると０になる。
  　$ \delta(p,x) $が隠れた変数となり、異次元で誤差をこの不確定性原理
  と表している。
  全ては素粒子方程式から生成される式達である。
   $$ (u + d) + c = e^{-f} - e^{f}, b + (w + s) = e^{-f} + e^{f} $$
   $$ \int \int{1 \over x^s}\mathrm{dvol} - \int \log x \mathrm{dvol}
   = {d \over df}\int\int{1 \over (y \log y)^{1 \over 2}}dy_m
   - \int C dx_m $$
   $$ = e^{-f} - e^{f} $$
   $$ {d \over df}F = {d \over df}\int \int{1 \over (x \log x)^2}dx_m 
   + {d \over df}\int\int{1 \over (y \log y)^{1 \over 2}}dy_m $$
   $$ = e^{f} + e^{-f} $$

     素粒子方程式は、２種類の素粒子と１種類の素粒子の組み合わせで、
   Jones多項式をペアーで構成されている。Rico level理論をベースに
   して、強い力と弱い力、電磁気力が、誤差をD-brane間を重力が
   gravitonとして伝わっているのが、洩れて、反重力をこの
   Weak electric theoryに加えると、全部の力が合わさって、
   統一場理論が完成される。
   これが、ヒッグス場とオイラーの定数を多様体を次元として、
   加群すると、ゼータ関数になる。これが、微分幾何の量子化とガンマ関数
   でもある。

   $$ ||ds^2|| = e^{-2 \pi T|\psi|}[\eta + \bar{h}(x)]dx^{\mu}dx^{\nu}
   + T^2 d^2 \psi $$
   
      素数分布論が、ゼータ関数を開集合として同型でもあり、ノルム空間、経路積分
   として、位相差が原子間力にもなっている。このエネルギー分布が、
   素数と素粒子方程式、宇宙と異次元のフェルミオンとして、一致している。
   このエネルギー分布が、複素多様体である。
   $$\square = 2 (\cos (i x \log x) + i \sin(i x \log x) $$ 


      宇宙と異次元が、双対性として、超対称性の素粒子を形成している。　


 $$ {\partial \over \partial f}F(x) = \int \int \text{cohom}D_k(x)[I_m]  $$
 Cohomology element is constructed with isotopy of D-brane,
 and this isotopy of symmetry structure is sheaf of manifold,
 more over this brane of element is double integrate of projection.
 $$ {\partial \over \partial f}F = {{}^t {\scalebox{1}[2]{{\variint}}}}
 \text{cohom}D_k(x)^{\ll p} $$
 And global partial defferential equation is integrate of cut in cohomolgy,
 and this step of defferential D-brane of sheap value is varint of equals,
 inverse of horizen of cute of section value is reverse of integrate of 
 operators. Three dimension of varint of manifolds in operator of sheaps,
 and this step of times of value is equal with D-brane of equations.
 Global differential equation and integrate equation is operator of 
 global scales of horizen cut and add of cut equation are routs.

 $$ \oint r dx_m = \mathcal{O}(x,y) $$

 $$ {{}^t {\scalebox{1}[2]{{\variiint}}}}_{{D(\chi,x)}} \text{Hom}[D^2 \psi]
 ^{\ll p} \cong  \text{vol}\left({V \over S}\right) $$

 $$ {\partial \over \partial f}F(x) = {{}^t \scalebox{1}[2]{\variint}} 
 \text{cohom}D_k(x)^{\ll p} $$

 $$ {d \over df}F = (F)^{f^{'}}, \int F dx_m = (F)^{f}  $$
 $$ {d \over df}\int \int{1 \over (x \log x)^2}dx_m =
 \left(\int \int{1 \over (x \log x)^2}dx_m \right)^{f^{'}}  $$
 $$ f^{'} = (2x (\log x + \log (x + 1))) $$
  $$  \int \int F dx_m = \left({1 \over (x \log x)^2}\right)^{
	  (\int 2x (\log x + \log (x + 1))dx)} $$ 
   $$= e^{-f} $$
  $$ {d \over df}F = \left({1 \over (x \log x)^2}\right)^{
	  (2x (\log x + \log (x + 1))) } $$  
  $$= e^{f} $$
  $$ \log (x \log x) \ge 2(y \log y)^{1 \over 2} $$

  $$ \log (x \log x) \ge 2(\sqrt{y \log y}) $$

  Global differential manifold is calucrate with x of x times in non
  entropy of defferential values, and out of range in differential formula
  be able to oneselves of global differential compute system.
  Global integrate manifold is also calucrate with step of resulted with
  loop of integrate systems. Global differential manifold is newton and 
  laiptitz formula of out of range in deprivate of compute systems.
  This system of calcurate formula is resulted for reverse of global equation
  each differential and intergrate in non entropy compute resulted values.

  $$ \pi(\chi,x) = [i\pi(\chi,x),f(x)] $$
  $$ \int{1 \over (x \log x)}dx = i \int{1 \over (x \log x)}dx^{N}  
  + {}^{N} i \int{1 \over (x \log x)}dx $$
  $$  \int{1 \over (x \log x)}dx = i \int{x \log x}dx 
  - \int{1 \over (x \log x)}dx $$
  $$ \int \int{1 \over (x \log x)^2}dx_m = i {1 \over 2}x^2 $$
  $$ \int \int{1 \over (x \log x)^2}dx_m = i \int \int_M dx_m $$
  $$ \le {1 \over 2}i + x^2 $$
  $$ E = - {1 \over 2}mv^2 + mc^2 $$
  $$ \lim_{x \to \infty} \int \int{1 \over (x \log x)^2}dx_m \ge {1 \over 2}i $$
  $$ {d \over df}\int \int{1 \over (x \log x)^2}dx_m = {1 \over 2}i $$
  $$ {d \over df}\int \int{1 \over (x \log x)^2}dx_m = 
  \left(\int \int {1 \over (x \log x)^2}dx_m \right)^{(f)^{'}}$$
  $$ (f)^{(f)^{'}} $$
  $$ = (f^f)^{'} $$
  $$ = e^{x \log x} $$

  $$ \Gamma = \int {e^{-x}x^{1-t}}dx, 
  {d \over df}F = e^f, \int F dx_m = e^{-f} $$
  $$ \int x^{1-t}dx = {d \over df}F, \int e^{-x}dx = \int F dx_m $$
  $$ e^{i \theta} = \cos \theta + i \sin \theta $$
  $$ x \bot y \to x \cdot y = \vec{0}, 
  {d \over df}\int \int{1 \over (x \log x)^2}dx_m = {1 \over 2}i $$

  $$ {d \over df}\int \int{1 \over (x \log x)^2}dx_m = 
  \left(\int \int {1 \over (x \log x)^2}dx_m \right)^{(f)^{'}}$$
  $$ {1 \over 2} + {1 \over 2}i = \vec{0}, {1 \over 2}\cdot{1 \over 2}i $$
  $$ A + B  =  \vec{0}, A \cdot B = {1 \over 4}i $$
  $$  \tan 90 \neq 0, ||ds^2|| = A + B $$
  Represented real pole is not only one of pole but also real pole of 
  $ {1 \over 2} $, 
  $\sin 0 = 0$ $$ e^{i \theta} = \cos \theta + i \sin \theta $$
  $\theta $
  this result equation is non-certain system concerned with particle, electric 
  particle and atoms in open set group with possiblity of surface values,
  and have abled to proof in this group of system. Universe and the other
  dimension of concerned of relationship values.
  $$ {d \over df}F = {d \over df} \int \int{1 \over (x \log x)^2}dx_m 
  + {d \over df} \int \int {1 \over (y \log y)^{1 \over 2}}dy_m
  = \Gamma(x,y) $$
  $$ {d \over d \gamma}\Gamma = (e^{-x}x^{1-t})^{\gamma^{'}} \to
  {d \over d \gamma}\Gamma = \Gamma^{(\gamma)^{'}}$$



    $$ \Gamma(s) = \int e^{-x} x^{s - 1}dx, \Gamma^{'}(s)
    = \int e^{-x} x^{s - 1}\log x dx $$
    $$ x^{s - 1} = e^{(s - 1)\log x} $$
    $$ {\partial \over \partial s}x^{s - 1} =
       {\partial \over \partial s}e^{-x}\cdot (e^{{s - 1}\log x}) 
       = e^{-x}\cdot \log x \cdot e^{(s - 1) \log x} $$
    $$ = e^{-x}\cdot \log x \cdot x^{s - 1} $$
    $$ {d \over d \gamma}\Gamma = \Gamma^{(\gamma)^{'}} $$
    $$ {d \over d \gamma}\Gamma(s) = \left( \int^{\infty}_{0}e^{-x}
    x^{s - 1}dx \right)^{\left(\int^{\infty}_{0}e^{-x}x^{s - 1}\log x dx
    \right)^{'}} $$
    $$ = \Gamma^{(\Gamma \int \log x dx)^{'}} $$
    $$ = e^{-x \log x} $$

    $$ \Gamma(s) = \int^{\infty}_{0} e^{-x}x^{s - 1}dx $$
    $$  \Gamma^{'}(s) = \int^{\infty}_{0} e^{-x}x^{s - 1} \log x dx $$
    $$ {d \over df}F = \int x^{s - 1}dx $$
    $$ \int F dx_m = \int e^{-x}dx $$
    $$ {d \over df}F = F^{(f)^{'}}, \int F dx_m = F^{(f)} $$
    Global differential manifold and global integrate manifold are
    reached with begin estimate of gamma and beta function leaded solved,
    Global differential manifold is emerged with gamma function of
    themselves of first value of global differential manifold of gamma
    function, this equation of quota in newton and dalarnverles of equation
    is also of same results equations. Zeta function also resulted with
    quantum group reach with quanto algebra equation.
        Included of diverse of gravity of influent in quantum physics,
    also quantum physics doesn't exist Gauss surface of theorem,
   
    $$ H \Psi = \bigoplus (i \hbar ^{\nabla})^{\oplus L} $$
    $$ = \bigoplus {{ H \Psi} \over \nabla L} $$
    $$ = e^{x \log x} = x^{(x)^{'}} $$
    however,
    $$ {d \over df}F = m(x) $$
    and gravity equation in being abled to existed with quantum physics,
    and this physics is also represented with quantum level of differential
    geometry.
    $$ {d \over dt}\psi(t) = \hbar $$
    $$ = {1 \over 2}i e^{i \Hat{H}} $$
    $$ (i \hbar)^{'} = (- e^{i \Hat{H}})^{'} $$
    $$ = - i e^{i \Hat{H}} $$
    $$ \psi(x) = e^{-i \Hat{H}t}, \bigoplus (i \hbar ^{\nabla})^{\oplus L} 
    = {{{1 \over 2} e^{i \Hat{H}}}^{(-i e^{i \Hat{H}})}} $$
    $$ = ({1 \over 2}f)^{-i f}$$
    $$ = ({1 \over 2})^{-i f} \cdot e^{-x \log x} \cdot (f)^{i} $$
    $$ = \int e^{- x} x^{t - 1}dx, {d \over d \gamma}\Gamma = e^{-x \log x}  $$
     
    $$ f = x, i = t, {1 \over 2} = a, \bigoplus a^{-t x}x^t [I_m]
     \cong \int e^{-x} x^{t - 1}dx $$
     Quantum level of differential geomerty is also resulted with
    Gamma function of global differential manifold of values, Heisenberg
    equation is alos resulted with global integrate manifold and this
    quantum level of differential geometry is manifold of deprive of formula.
    $$ |\psi(t)\rangle_s = e^{-i\Hat{H} t}|\Psi\rangle_{H},
    \Hat{A}_s = \Hat{A}_H(0) $$
    $$ |\Psi(t)\rangle_s \to {d \over dt} $$
    $$ i {d \over dt}|\psi (t) \rangle_s = \Hat{H}|\psi(t)\rangle_s $$
    $$ \langle\Hat{A}(t)\rangle = \langle\Psi(t)|\Hat{A}(0)|\Psi(t)\rangle  $$
    $$ {d \over dt}\Hat{A} = {1 \over i}[\Hat{A},H] $$
    $$ \Hat{A}(t) = e^{i\Hat{H}t}\Hat{A}(0)e^{-i\Hat{H}t} $$
    $$ \lim_{\theta \to 0} \begin{pmatrix} \sin \theta \\
    \cos \theta \end{pmatrix} \begin{pmatrix} \theta & 1 \\
    1 & \theta \end{pmatrix} \begin{pmatrix} \cos \theta \\
    \sin \theta \end{pmatrix} = \begin{pmatrix} 1 & 0 \\
    0 & - 1 \end{pmatrix} $$
    $$ f^{-1}(x) x f(x) = I^{'}_m, I^{'}_m = [1,0] \times [0,1] $$

    $$ {x + y} \ge \sqrt{x y} $$
    $$ {x^{{1 \over 2} + iy} \over e^{x \log x}} = 1 $$
    $$ \mathcal{O}(x) = \nabla_i \nabla_j \int e^{{2 \over m}\sin \theta 
    \cos \theta } \times {{N \mathrm{mod}(e^{x \log x})} \over {O(x)(x +
    \Delta |f|^2)^{1 \over 2}}} $$
    $$ x \Gamma(x) = 2 \int|\sin 2 \theta|^2 d \theta, \mathcal{O}(x)
    = m(x)[D^2 \psi] $$
    $$ i^2 = (0,1) \cdot (0,1), |a||b|\cos \theta = - 1 $$
    $$ E = \mathrm{div}(E, E_1) $$
    $$ \left({\{f,g\}} \over [f,g] \right) = i^2, E = mc^2, I^{'} = i^2 $$
     Gamma function of global differential equation is first of differential
    function deprivate with ordinary differential equations, more also
    universe and other dimensiion of Higgs field emerge with zeta function,
    more spectrum focus is three dimension of manifold of singularity or pair
    theorem, and universe and other dimension of each value equals.
    This reverse of circle function of manifolds is also Higgs fields of 
    dense of entropy and result equation. And differential structure of quantum
    level in gravity is poled with zeta function and quantum group in high 
    result values. Zeta funtion and quantum group is Global differential
    manifold of result in values.
    $$ {d \over d \gamma} \Gamma = m(x) $$
    $$ = e^{- x \log x} $$
    $$ \sin i x = {{e^{-x} + e^{x}} \over 2i} $$
    $$ {d \over df}F = m(x) = e^{f} + e^{-f} $$
    $$ = 2 i \sin (i {x \log x}) $$

    Hyper circle function is constructed with sheap of manifold in 
    Beta function of reverse system emerged from Gamma fucntion 
    and reverse of circle function of entropy value,
    and this reverse of hyper circle function of beta system
    is universe and the other dimension of zeta function,
    more over spectrum focus is this beta function of reverse in
    circle function is quantum level of differential structure oneselves,
    and this quantum level is also gamma function themselves.
    This gamma fucntion is more also D-brane and anti-D-brane built 
    with supplicant values. Higgs fileds is more also pair of universe and
    other dimension each other value elements. 


    My son in acasicrecord demonstrate with being have with quantum level
    of differential structure equation.
    I have with same equation already from being on 
    my father and my doctors of professors give me hints with 
    this qlds equation explanade with global differential manifold.

    $$ H\Psi = \bigoplus {(i \hbar^{\nabla})}^{\oplus L}  (1)$$
    $$ = i\hbar\psi $$
    $$ {d \over df}F = m(x) (2) $$
    $$ = {d \over df}\int \int{1 \over (x \log x)^2}dx_m
    + {d \over df}\int \int{1 \over (y \log y)^{1 \over 2}}dy_m  (3)$$
    $$  \bigoplus {(i \hbar^{\nabla})}^{\oplus L} = \int e^{-x}x^{1 - t}dx_m 
    (4)$$
    $$ = {d \over d \gamma}\Gamma = \int \Gamma(\gamma)^{'}dx_m (5)$$
    $$ = e^f + e^{-f} (6)$$
    $$ {d \over d \gamma}\Gamma^{-1} = e^{f} - e^{-f} (7)$$
    (1) is my son in acasicrecord with his idea, (2),(3) are myself equation.
    (4),(5) are my son and me with tagge of ideas equation.
    (6),(7) are Takashima Aya and me, and my son in acasicrecode with tolio
    equations.
    These references from Grisha professor and Takeuchi Kaoru.
    I read with this references being hint from these professors.
 
      [reference Lisa Randall, Toshihide Masuoka, Makoto Kobayashi]


 
       
\end{document}
